%% ============================================================
%%  ਹਵਨ - ਇੱਕ ਪੁਰਾਤਨ ਮਰਿਯਾਦਾ  (Nanak Niwas)
%%  Built from base_template.  Compile with:  make pdf
%% ============================================================
\DocumentMetadata{
  lang=pa,
  pdfversion=2.0,
  tagging=on,
  tagging-setup={interwordspace=true}
}
\documentclass[12pt,a5paper,twoside]{book}  % use oneside to avoid blank pages

\usepackage[a5paper, top=1.5cm, bottom=2cm, inner=2cm, outer=1.2cm]{geometry}

\usepackage[table]{xcolor}
\usepackage{pagecolor}
\usepackage{graphicx}
\usepackage{fontspec}
\graphicspath{{images/}{images/gallery/}}
\usepackage{polyglossia}
\setmainlanguage{punjabi}
\usepackage{enumitem}
\usepackage{ragged2e}
\usepackage{longtable}
\usepackage{tikz}     % For drawing the page border
\usepackage{eso-pic}  % To repeat the border on every page
\setlist{topsep=0.3em, itemsep=0.3em}

\newfontfamily\punjabifont[Script=Gurmukhi,
  Path=fonts/, Renderer=HarfBuzz,
  BoldFont=NotoSerifGurmukhi-SemiBold.ttf
]{NotoSerifGurmukhi-Regular.ttf}

\usepackage{fancyhdr}
\pagestyle{fancy}
\fancyhf{}
\fancyfoot[C]{\thepage}
\renewcommand{\headrulewidth}{0pt}
\renewcommand{\footrulewidth}{0pt}

\definecolor{nanak_niwas_midnight_violet}{HTML}{241236}
\definecolor{nanak_niwas_thistle}{HTML}{D8CCE6}
\definecolor{nanak_niwas_soft_peach}{HTML}{F2D6A2}
\definecolor{nanak_niwas_metalic_gold}{HTML}{D4AF37}
\newcommand{\currentBorderColor}{nanak_niwas_metalic_gold}

\usepackage{setspace}
\setstretch{1.25}
\setlength{\parindent}{0pt}
\setlength{\parskip}{0.5em}

% Heading styles (set by hand: titlesec is not compatible with tagged PDF).
% Applied at the very end of begin{document} so that the tagging code
% cannot overwrite them.
\makeatletter
\AddToHook{begindocument/end}{%
  \RenewDocumentCommand\chapter{s o m}{%
    \clearpage
    \IfBooleanTF{#1}{%
      {\centering\Huge\bfseries\punjabifont #3\par\nobreak\vskip 1em}%
      \@afterheading
    }{\@latex@error{Only \string\chapter* is supported here}\@ehc}%
  }%
  \renewcommand\section{\@startsection{section}{1}{0pt}{1.2em}{0.6em}{\punjabifont\Large\bfseries}}%
  \renewcommand\subsection{\@startsection{subsection}{2}{0pt}{1em}{0.5em}{\punjabifont\large\bfseries}}%
}
\makeatother

% Page border on every page
\AddToShipoutPictureBG{%
  \begin{tikzpicture}[remember picture, overlay]
    \draw[line width=3pt, color=\currentBorderColor]
      ([shift={(0.7cm, 0.7cm)}]current page.south west)
      rectangle
      ([shift={(-0.7cm, -0.7cm)}]current page.north east);
  \end{tikzpicture}%
}

% Table of contents title as in the source text
\addto\captionspunjabi{\renewcommand{\contentsname}{ਤਤਕਰਾ}}

%% ---------------- BOOK DETAILS ----------------
\newcommand{\booktitle}{ਹਵਨ}
\newcommand{\booksubtitle}{ਇੱਕ ਪੁਰਾਤਨ ਮਰਿਯਾਦਾ \\[0.6em] {\Large ਭਾਈ ਮਨਪ੍ਰੀਤ ਸਿੰਘ}}
\newcommand{\coverimage}{cover.jpg}
\newcommand{\logoimage}{nanak_nivas.png}

%% ---------------- Helpers ----------------
% Latin-script text (the Gurmukhi font has no Latin letters)
\DeclareRobustCommand{\eng}[1]{{\rmfamily #1}}

% Full-page gallery photo:  \galleryphoto{file}{caption}
\newcommand{\galleryphoto}[2]{%
  \begin{figure}[p]
    \thispagestyle{empty}
    \centering
    \includegraphics[width=\textwidth,height=0.85\textheight,keepaspectratio,alt={#2}]{#1}
    \vspace{1em}
    \caption{#2}
  \end{figure}
  \clearpage
}

\begin{document}

{\punjabifont

%% ---------------- COVER ----------------
\begin{titlepage}
    \newpagecolor{nanak_niwas_midnight_violet}
    \color{nanak_niwas_metalic_gold}
    \thispagestyle{empty}

    \newgeometry{top=2cm, bottom=2cm, left=1.5cm, right=1.5cm}
    \centering

    {\large \bfseries ੴ ਸਤਿਗੁਰ ਪ੍ਰਸਾਦ \par}

    \vspace*{0.3cm}

    \fboxsep=0pt
    \fboxrule=0.5pt
    \fbox{\includegraphics[width=0.72\textwidth,alt={ਕਵਰ ਚਿੱਤਰ}]{\coverimage}}

    \vfill

    {\Huge \bfseries \booktitle \\
    \booksubtitle \par}
    \vfill
    \fbox{\includegraphics[width=0.2\textwidth,alt={ਨਾਨਕ ਨਿਵਾਸ ਲੋਗੋ}]{\logoimage}}
    \restoregeometry
\end{titlepage}
\restorepagecolor
\renewcommand{\currentBorderColor}{nanak_niwas_midnight_violet}

%% ---------------- FRONT MATTER ----------------
\null
\thispagestyle{empty}
\newpage
\tableofcontents

\chapter*{ਭੂਮਿਕਾ}
\addcontentsline{toc}{chapter}{ਭੂਮਿਕਾ}
\section*{ਸਮਰਪਿਤ}
\subsection*{ਧੰਨ ਮਾਤਾ ਸਾਹਿਬ ਦੇਵਾਂ ਜੀ (ਖਾਲਸਾ ਪੰਥ ਦੇ ਪੂਜਨੀਕ ਮਾਤਾ ਜੀ)}
ਮੈਂ ਮਾਤਾ ਸਾਹਿਬ ਦੇਵਾਂ ਜੀ ਦੇ ਚਰਨਾ ਕਮਲਾਂ ਵਿਚ ਨਿਮਰਤਾ ਸਹਿਤ ਪ੍ਰਣਾਮ ਕਰਦਾ ਹਾਂ ਅਤੇ ਇਸ ਪੁਸਤਕ ਨੂੰ ਸੇਵਾ ਵਜੋਂ ਭੇਟ ਕਰਦਾ ਹਾਂ। ਬ੍ਰਹਮ ਸ਼ਕਤੀ ਦੇ ਬਗੈਰ ਮੈਂ ਹਿੱਲਣ ਦੇ ਵੀ ਯੋਗ ਨਹੀਂ ਸੀ। ਇਹ ਉਹ ਸ਼ਕਤੀ ਹੈ ਜਿਹਨਾਂ ਦੇ ਸਹਿਯੋਗ ਦੇ ਨਾਲ ਇਹ ਕਿਤਾਬ ਲਿਖੀ ਗਈ ਹੈ। ਮੈਂ ਫੇਰ ਆਪਣਾ ਸੀਸ ਧੰਨ ਗੁਰੂ ਨਾਨਕ ਦੇਵ ਸਾਹਿਬ ਜੀ ਨਿਰੰਕਾਰ ਨੂੰ ਭੇਟਾ ਕਰਦਾ ਹਾਂ ਜੋ ਮੇਰੇ ਮਾਰਗ ਦਰਸ਼ਕ ਹਨ। ਫੇਰ ਮੈਂ ਆਪਣੀ ਫ਼ਤਹਿ ਧੰਨ ਬਾਬਾ ਸ੍ਰੀ ਚੰਦ ਸਾਹਿਬ ਜੀ (ਧੰਨ ਗੁਰੂ ਨਾਨਕ ਦੇਵ ਸਾਹਿਬ ਜੀ ਦੇ ਸਾਹਿਬਜ਼ਾਦੇ) ਨੂੰ ਭੇਟ ਕਰਦਾ ਹਾਂ ਜਿਹਨਾਂ ਨੇ ਮੇਰੇ ਅੰਦਰ ਪੁਰਾਤਨ ਮਰਿਯਾਦਾ ਨੂੰ ਜਿਉਂਦਾ ਰੱਖਿਆ।

ਧੰਨ ਬਾਬਾ ਸ੍ਰੀ ਚੰਦ ਜੀ ਮਹਾਰਾਜ, ਧੰਨ ਬਾਬਾ ਬਲਵੰਤ ਸਿੰਘ ਜੀ (ਸਿਹੜਾ ਸਾਹਿਬ), ਧੰਨ ਬਾਬਾ ਵਿਰਸਾ ਸਿੰਘ ਜੀ (ਗੋਬਿੰਦ ਸਦਨ) ਜਿਨ੍ਹਾਂ ਦੇ ਦਾਸ ਦੇ ਸਿਰ ਤੇ ਹੱਥ ਰੱਖ ਕੇ ਇਹ ਗ੍ਰੰਥ ਦੀ ਸੇਵਾ ਕਰਾਈ ਹੈ।

ਮੈਂ ਸਾਧ ਸੰਗਤ (ਸਰਦਾਰ ਚਰਨਪ੍ਰੀਤ ਸਿੰਘ, ਬੀਬੀ ਗੁਰਿੰਦਰ ਕੌਰ, ਬੀਬੀ ਹਰਸ਼ੀਨ ਕੌਰ, ਬੀਬੀ ਜਸਲੀਨ ਕੌਰ) ਅਤੇ ਆਪਣੇ ਮਾਤਾ ਪਿਤਾ ਜੀ ਦਾ ਧੰਨਵਾਦ ਕਰਦਾ ਹਾਂ ਜਿਹਨਾਂ ਦੇ ਸਹਿਯੋਗ ਨਾਲ ਇਸ ਪੁਸਤਕ ਵਿਚ ਮੌਜੂਦ ਜਾਣਕਾਰੀ ਇਕੱਠੀ ਕੀਤੀ ਗਈ। ਜੇ ਸਰੋਤਿਆਂ ਨੂੰ ਇਸ ਕੰਮ ਵਿੱਚ ਕੋਈ ਕਮੀਆਂ ਨਜ਼ਰ ਆਉਂਦੀਆਂ ਹਨ ਤਾਂ ਉਹ ਸਿਰਫ ਮੇਰੀਆਂ ਹਨ, ਕਿਰਪਾ ਕਰਕੇ ਮੈਨੂੰ ਮੁਆਫ਼ ਕਰਨਾ ਅਤੇ ਦਾਸ ਨੂੰ ਆਪਣੇ ਵਿਚਾਰ ਜ਼ਰੂਰ ਭੇਜਣਾ।
\newpage
ਇਹ ਸੇਵਾ ਨਾਨਕ ਨਿਵਾਸ ਵੱਲੋਂ ਨਿਮਾਣੇ ਰੂਪ ਵਿੱਚ ਭੇਟ ਕੀਤੀ ਗਈ ਹੈ। ਇਸ ਵਿਸ਼ੇਸ਼ ਦਸਤਾਵੇਜ਼ ਦੀ ਸੇਵਾ ਹੇਠ ਲਿਖੇ ਸੇਵਾਦਾਰ ਵੱਲੋਂ ਪੂਰੀ ਸ਼ਰਧਾ ਨਾਲ ਕੀਤੀ ਗਈ ਹੈ:
\begin{itemize}
    \item ਭਾਈ ਮਨਪ੍ਰੀਤ ਸਿੰਘ
\end{itemize}
\clearpage

%% ---------------- NANAK NIWAS ----------------
\chapter*{ਨਾਨਕ ਨਿਵਾਸ}
\addcontentsline{toc}{chapter}{ਨਾਨਕ ਨਿਵਾਸ}
ਨਾਨਕ ਨਿਵਾਸ ਇੱਕ ਗੁਰਮਤਿ ਵਿਦਿਆਲਿਆ ਹੈ ਜਿਸਦਾ ਮਕਸਦ ਗੁਰੂ ਨਾਨਕ ਦੇਵ ਜੀ ਦੀ ਅਦਵੈਤ ਅਤੇ ਪ੍ਰੇਮ ਦੀ ਸਿੱਖਿਆ ਨੂੰ ਸਾਂਝਾ ਕਰਨਾ ਹੈ, ਅਤੇ ਖੋਜੀਆਂ ਨੂੰ ਆਪਣੇ ਅਸਲ ਸਰੂਪ ਨੂੰ ਸਮਝਣ ਵਿੱਚ ਮਦਦ ਕਰਨਾ ਹੈ।
\\
ਬ੍ਰਹਮ ਵਿਦਿਆ ਦੀਆਂ ਕਲਾਸਾਂ ਹਰ ਹਫਤੇ ਆਨਲਾਈਨ ਪੰਜਾਬੀ, ਹਿੰਦੀ ਅਤੇ ਅੰਗ੍ਰੇਜ਼ੀ ਵਿੱਚ ਹੁੰਦੀਆਂ ਹਨ, ਜਿਥੇ ਅਸੀਂ ਪੁਰਾਤਨ ਗ੍ਰੰਥਾਂ ਨੂੰ ਗੁਰਮਤਿ ਦੀ ਨਜ਼ਰ ਨਾਲ ਸਮਝਦੇ ਹਾਂ ਅਤੇ ਇਹ ਵੇਖਦੇ ਹਾਂ ਕਿ ਇਹ ਗਿਆਨ ਅਸੀਂ ਆਪਣੀ ਰੋਜ਼ਾਨਾ ਜ਼ਿੰਦਗੀ ਵਿੱਚ ਕਿਵੇਂ ਲਾਗੂ ਕਰ ਸਕਦੇ ਹਾਂ। ਇਸਦੇ ਨਾਲ ਨਾਲ, ਅਸੀਂ ਕਾਵਿ ਵਿਦਿਆ ਦੀਆਂ ਆਨਲਾਈਨ ਕਲਾਸਾਂ ਵੀ ਕਰਵਾਉਂਦੇ ਹਾਂ ਅਤੇ ਮੇਲਬਰਨ ਵਿੱਚ ਹਰ ਮਹੀਨੇ ਸੰਗਤ ਦੇ ਸੈਸ਼ਨ ਰੱਖਦੇ ਹਾਂ, ਜਿਥੇ ਸਿੱਖਣ, ਜੁੜਨ ਅਤੇ ਇਕੱਠੇ ਵਧਣ ਦਾ ਮੌਕਾ ਮਿਲਦਾ ਹੈ।
\\
ਨਾਨਕ ਨਿਵਾਸ ਹਰ ਉਸ ਵਿਅਕਤੀ ਲਈ ਹੈ ਜੋ ਸੱਚ ਦੀ ਖੋਜ ਵਿੱਚ ਹੈ, ਕਿਸੇ ਇੱਕ ਸਮਾਜ ਜਾਂ ਸੰਪ੍ਰਦਾ ਤੱਕ ਸੀਮਤ ਨਹੀਂ। ਸਾਲ ਭਰ ਅਸੀਂ ਵੱਖ-ਵੱਖ ਸਮਾਗਮ ਵੀ ਕਰਦੇ ਹਾਂ ਤਾਂ ਕੀ ਸੰਗਤ ਨਾਲ ਹੋਰ ਜੁੜ ਸਕੀਏ ਅਤੇ ਇਹ ਸੰਦੇਸ਼ ਸਾਂਝਾ ਕਰ ਸਕੀਏ, ਨਾਲ ਹੀ ਸੰਗਤਾਂ ਨੂੰ ਆਪਣੀ ਖੁਦ ਦੀ ਖੋਜ ਕਰਨ ਲਈ ਪ੍ਰੇਰਿਤ ਕਰਦੇ ਹਾਂ ਤਾਂ ਕੀ ਉਹ ਸਿਰਫ਼ ਮੰਨਣ ਦੀ ਬਜਾਏ ਇਸਨੂੰ ਆਪਣੀ ਜ਼ਿੰਦਗੀ ਵਿੱਚ ਜੀ ਸਕਣ।
\begin{figure}[h]
      \centering
      \includegraphics[width=4cm,keepaspectratio,alt={ਨਾਨਕ ਨਿਵਾਸ ਲੋਗੋ}]{nanak_nivas.png}
    \end{figure}

%% ================= BOOK CHAPTERS =================

%% ---------------- 1 ----------------
\clearpage
\chapter*{\eng{1.} ਹਵਨ ਕੀ ਹੈ?}
\addcontentsline{toc}{chapter}{\eng{1.} ਹਵਨ ਕੀ ਹੈ?}
ਇਹ ਪੁਸਤਕ ਖਾਲਸਾ ਪੰਥ ਦੀ ਪੁਰਾਤਨ ਹਵਨ ਪਰੰਪਰਾ ਨੂੰ ਸੁਰਜੀਤ ਕਰਨ ਅਤੇ ਉਸ ਵਿਚ ਰੌਸ਼ਨੀ ਲਿਆਉਣ ਦਾ ਇੱਕ ਉਪਰਾਲਾ ਹੈ, ਜੋ ਕਿ ਦਲ ਪੰਥ (ਬੁੱਢਾ ਦਲ ਅਤੇ ਤਰਨਾ ਦਲ) ਅਤੇ ਹੋਰ ਨਿਹੰਗ ਸਿੰਘ ਸੰਪਰਦਾਵਾਂ ਵਿਚ ਅਜੇ ਵੀ ਚਲਾਈ ਜਾਂਦੀ ਹੈ। ਸਾਲਾਂ ਦੌਰਾਨ, ਖਾਲਸੇ ਵਿੱਚ ਸੰਨ ੧੮੭੦ ਵਿਚ ਸ਼ੁਰੂ ਹੋਈ ਸਿੰਘ ਸਭਾ ਲਹਿਰ ਦੇ ਨੁਕਸਾਨ ਵਿੱਚ ਇੱਕ ਸੀ। ਕਿਉਂਕਿ ਕੁਝ ਅਭਿਆਸ ਜੋ ਪੁਰਾਤਨ ਸਿੱਖਾਂ ਦੁਆਰਾ ਕੀਤੇ ਜਾਂਦੇ ਸੀ ਉਹਨਾਂ ਨੂੰ ਨਿਰਉਤਸ਼ਾਹਿਤ ਕੀਤਾ ਅਤੇ ਹਿੰਦੂਆਂ ਦੀ ਪਰੰਪਰਾ ਅਤੇ ਕੁਝ ਮੁਸਲਮਾਨਾਂ ਦੀ ਪਰੰਪਰਾ ਕਰਕੇ ਬਰਖਾਸਤ ਕੀਤਾ ਗਿਆ। ਇਹ ਪੁਸਤਕ ਇਸ ਗੱਲ ਦੀ ਡੂੰਘਾਈ ਨਾਲ ਖੋਜ ਨਹੀਂ ਕਰੇਗੀ ਕਿ ਸਿੰਘ ਸਭਾ ਲਹਿਰ ਦੇ ਕਾਰਨ ਕੀ ਅਤੇ ਕਿਵੇਂ ਵਾਪਰਿਆ, ਪਰ ਹਵਨ ਦੀ ਪਰੰਪਰਾ ਅਤੇ ਗੁਰਮਤਿ ਅਤੇ ਵਿਗਿਆਨਿਕ ਪੱਖੋਂ ਇਸ ਦੀ ਮਹੱਤਤਾ ਤੇ ਜ਼ੋਰ ਦਿੱਤਾ ਜਾਵੇਗਾ। ਜਿਵੇਂ ਕਿ ਹਵਨ ਨੂੰ ਹਿੰਦੂਆਂ ਵੱਲੋਂ ਵੱਡੇ ਪੱਧਰ ਤੇ ਅਭਿਆਸ ਕਰਨ ਲਈ ਜਾਣਿਆ ਜਾਂਦਾ ਹੈ, ਉਹਨਾਂ ਦੀ ਪਰੰਪਰਾ ਦਾ ਕੁਝ ਹਵਾਲਾ ਵੀ ਤੁਲਨਾ ਦੇ ਅਧਾਰ ਤੇ ਸਿੱਟਾ ਕੱਢਣ ਲਈ ਬਣਾਇਆ ਜਾਵੇਗਾ।

ਉਹਨਾਂ ਲਈ, ਜਿਹਨਾਂ ਕੋਲ ਖੁੱਲ੍ਹੇ ਦਿਮਾਗ ਅਤੇ ਸਿੱਖਣ ਦੀ ਇੱਛਾ ਹੈ, ਇਹ ਕਿਤਾਬ ਇੱਕ ਝਲਕ ਪ੍ਰਦਾਨ ਕਰਨ ਵਿਚ ਸਹਾਈ ਹੋਵੇਗੀ। ਜੇ ਤੁਸੀਂ ਉਹਨਾਂ ਵਿੱਚੋਂ ਹੋ ਜੋ ਕਿਤਾਬਾਂ ਅਤੇ ਲੇਖਾਂ ਨੂੰ ਇਤਰਾਜ ਉਠਾਉਣ ਅਤੇ ਪੁਰਾਤਨ ਪਰੰਪਰਾਵਾਂ ਦੀ ਨਿੰਦਿਆ ਕਰਨ ਲਈ ਇੱਕ ਸਾਧਨ ਵਜੋਂ ਵਰਤਦੇ ਹਨ ਤਾਂ ਮੈਂ ਉਹਨਾਂ ਪਾਠਕਾਂ ਨੂੰ ਨਿਮਰਤਾ ਸਹਿਤ ਬੇਨਤੀ ਕਰਦਾ ਹਾਂ ਕਿ ਉਹ ਅੱਗੇ ਨਾ ਪੜ੍ਹਨ ਕਿਉਂਕਿ ਮੈਂ ਤੁਹਾਡੇ ਅੰਦਰ ਕੋਈ ਗੁੱਸਾ ਨਹੀਂ ਪੈਦਾ ਕਰਨਾ ਚਾਹੁੰਦਾ ਪੂਰਵ ਧਾਰਨਾ ਅਤੇ ਮਜਬੂਰੀ ਨਾਲ ਜੜ੍ਹਾਂ ਵਾਲੀਆਂ ਝੂਠੀਆਂ ਧਾਰਨਾਵਾਂ (ਜੋ ਹਉਮੈ ਦੇ ਮਹਾਨ ਕੰਮ ਹਨ) ਤੋਂ ਪੈਦਾ ਹੁੰਦਾ ਹੈ।

ਜੇ ਤੁਹਾਡੇ ਕੋਲ ਕੋਈ ਸੁਝਾਅ ਜਾਂ ਸਵਾਲ ਹਨ, ਤਾਂ ਮੇਰੀ ਈ-ਮੇਲ ਤੇ ਮੇਰੇ ਤੱਕ ਪਹੁੰਚਾਉਣ ਲਈ ਤੁਹਾਡਾ ਸਵਾਗਤ ਹੈ: \eng{mannvinder@gmail.com}

%% ---------------- 2 ----------------
\clearpage
\chapter*{\eng{2.} ਅੱਗ ਕੀ ਹੈ?}
\addcontentsline{toc}{chapter}{\eng{2.} ਅੱਗ ਕੀ ਹੈ?}
“ਅਦਭੁਤ ਹੈ ਉਹ ਅੱਗ ਜੋ ਅਚੰਭੇ ਕਰਦੀ ਹੈ” (ਵਾਰ ਆਸਾ, ਮ:੧, ੪੬੪)। ਅਗਨੀ ਨੂੰ ਸੁਅੰਤਰ ਦੇਵਤਾ ਕਿਹਾ ਜਾਂਦਾ ਹੈ (ਵਾਰ ਆਸਾ, ਮ:੧, ੪੭੩)। ਇਸ ਦੇ ਹੋਰ ਨਾਮ ਪਾਵਕ ਅਤੇ ਤੇਜ ਵੀ ਹਨ। ਇਹ ਹਵਾ (ਪਵਨ ਦੇਵਤਾ) ਤੋਂ ਵਿਕਸਿਤ ਹੈ ਤੇ ਪ੍ਰਭੂ ਦੇ ਡਰ ਦੇ ਅਧੀਨ ਉਸ ਵਾਂਗ ਮਿਹਨਤ ਕਰਦੀ ਹੈ। ਜਦੋਂ ਸੂਰਜ ਅਸਮਾਨ ਵਿੱਚ ਚਮਕਦਾ ਹੈ, ਧਰਤੀ ਭੁੰਨ ਜਾਂਦੀ ਹੈ, ਅਤੇ ਉਸ ਦਾ ਸਰੀਰ ਸੂਰਜ ਦੀ ਅੱਗ ਨਾਲ ਸੁੱਕ ਜਾਂਦਾ ਹੈ (ਬਾਰਾਮਾਹਾ ਤੁਖਾਰੀ ਮ:੧, ੧੧੦੮)। ਸੂਰਜ ਤਾਪ (ਅਗਨੀ) ਦਾ ਮੁੱਖ ਸ੍ਰੋਤ ਹੈ।

ਅੱਗ ਲੱਕੜੀ ਵਿਚ ਲੁਕੀ ਹੋਈ ਹੈ। ਧੰਨ ਗੁਰੂ ਰਾਮਦਾਸ ਜੀ ਫਰਮਾਉਂਦੇ ਹਨ: ਅੱਗ ਲੱਕੜੀ ਵਿਚ ਬੰਦ ਹੈ, ਇਸ ਨੂੰ ਕੇਵਲ ਉਹੀ ਮਨੁੱਖ ਪਰਗਟ ਕਰਦਾ, ਜੋ ਵਿਧੀ ਜਾਣਦਾ ਹੈ (ਕਲਿਆਨ ਮ:੪, ਅੰਗ ੧੩੨੩)। ਜਿਵੇਂ ਗਰਭ ਦੀ ਅੱਗ ਸਰੀਰ ਅੰਦਰ ਹੈ, ਮਾਇਆ ਦੀ ਅੱਗ (ਸਰੀਰ ਵਿੱਚ) ਬਾਹਰ ਹੈ (ਰਾਮਕਲੀ ਮ: ੩, ਆਨੰਦ, ਅੰਗ ੯੨੧)।

ਇੱਛਾ ਵੀ ਇਕ ਕਿਸਮ ਦੀ ਅੱਗ ਹੈ, ਜਿਸ ਨੂੰ ਪ੍ਰਭੂ ਦੇ ਨਾਮ ਦੁਆਰਾ ਬੁਝਾਇਆ ਜਾ ਸਕਦਾ ਹੈ। ਜੰਗਲ ਦੀ ਅੱਗ (ਦਵਾ ਅਗਣ) ਦਾ ਵੀ ਜ਼ਿਕਰ ਹੈ, ਜੋ ਜੰਗਲ ਦੇ ਪੌਦਿਆਂ, ਰੁੱਖਾਂ ਦੇ ਨਾਲ ਨਾਲ ਕੀੜੇ ਮਕੌੜੇ, ਪੰਛੀਆਂ ਅਤੇ ਜਾਨਵਰਾਂ ਨੂੰ ਤਬਾਹ ਕਰ ਦਿੰਦੀ ਹੈ (ਅੰਗ: ੧੨੦, ੭੩੨, ੧੪੧੬)। ਅੱਗ ਦੀ ਇੱਕ ਛੋਟੀ ਜਿਹੀ ਚਿੰਗਾਰੀ ਬਾਲਣ ਦੇ ਵੱਡੇ ਢੇਰ ਨੂੰ ਸਾੜ੍ਹਨ ਲਈ ਕਾਫੀ ਹੈ (ਸੋਰਠ ਮ: ੧, ਅੰਗ ੬੩੭)। ਅੱਗ ਵੀ ਇੱਕ ਮਹਾਨ ਸ਼ੁੱਧ ਕਰਨ ਵਾਲੀ ਚੀਜ਼ ਹੈ। ਇਹ ਧਾਤ ਨੂੰ ਸ਼ੁੱਧ ਕਰਦੀ ਹੈ (ਸਲੋਕ ਫ਼ਰੀਦ, ਅੰਗ ੧੩੮੦)। ਨਰਕ ਦੀ ਅੱਗ ਪਾਪੀਆਂ ਲਈ ਹੈ। ਪਾਪੀ ਦਾ ਨਿਵਾਸ ਅੱਗ ਵਿੱਚ ਹੈ (ਭਗਤ ਨਾਮਦੇਵ, ਅੰਗ ੧੧੬੫) (ਮਾਰੂ ਮ: \eng{1}, ਅੰਗ ੧੦੩੩)। ਪਾਣੀ ਤੋਂ ਹੀ ਤਾਪ (ਅਗਨੀ) ਅਤੇ ਬਿਜਲੀ ਪੈਦਾ ਹੁੰਦੀ ਹੈ।

ਅਸੀਂ ਅਗਨੀ, ਅੱਗ ਦੇ ਦੇਵਤੇ ਬਾਰੇ ਇੱਕ ਹਵਾਲਾ ਨਿੱਜੀ ਤੌਰ ਤੇ ਕਰਦੇ ਹਾਂ, ਜਿਵੇਂ ਕਿ ਪੀ.ਥੌਮਸ ਦੁਆਰਾ (ਐਪਿਕਸ, ਮਿੱਥਸ ਐਂਡ ਲੀਜੰਡਸ, ਅੰਗ ੨੮) ਵਿੱਚ ਦਿੱਤਾ ਗਿਆ ਹੈ। ਇਹ ਇਸ ਵੈਦਿਕ ਦੇਵਤੇ ਦੀ ਮਹੱਤਤਾ ਨੂੰ ਦਰਸਾਉਂਦਾ ਹੈ: “ਅਗਨੀ ਮਨੁੱਖਾਂ ਦਾ ਸੁਆਮੀ, ਰੱਖਿਅਕ, ਰਾਜਾ ਹੈ। ਉਹ ਕਿਸੇ ਮਨੁੱਖ ਨੂੰ ਤੁੱਛ ਨਹੀਂ ਸਮਝਦਾ, ਉਹ ਹਰ ਪਰਿਵਾਰ ਵਿੱਚ ਰਹਿੰਦਾ ਹੈ। ਇਸ ਲਈ ਉਹਨੂੰ ਦੇਵਤਿਆਂ ਅਤੇ ਮਨੁੱਖਾਂ ਵਿਚਕਾਰ ਵਿਚੋਲੇ ਵਜੋਂ ਮੰਨਿਆ ਜਾਂਦਾ ਹੈ, ਅਤੇ ਉਹਨਾਂ ਦੀਆਂ ਕਾਰਵਾਈਆਂ ਦੇ ਗਵਾਹ ਵਜੋਂ। ਇਸ ਲਈ ਅੱਜ ਤਕ ਉਹਦੀ ਪੂਜਾ ਕੀਤੀ ਜਾਂਦੀ ਹੈ ਅਤੇ ਵਿਆਹ, ਮੌਤ ਦੇ ਸਾਰੇ ਧਾਰਮਿਕ ਮੌਕਿਆਂ ਤੇ ਉਸ ਦਾ ਅਸ਼ੀਰਵਾਦ ਮੰਗਿਆ ਜਾਂਦਾ ਹੈ।”

ਇਹਨਾਂ ਪੁਰਾਣੇ ਭਜਨਾਂ ਵਿਚ, ਅਗਨੀ ਨੂੰ ਲੱਕੜੀ ਦੇ ਦੋ ਟੁਕੜਿਆਂ ਦੇ ਵਿੱਚ ਨਿਵਾਸ ਕਰਨ ਬਾਰੇ ਕਿਹਾ ਗਿਆ ਹੈ, ਜਿਸ ਨੂੰ ਰਗੜ੍ਹਨ ਨਾਲ ਅੱਗ ਪੈਦਾ ਹੁੰਦੀ ਹੈ, ਅਤੇ ਇਹ ਇੱਕ ਕਮਾਲ ਦੀ ਗੱਲ ਹੈ ਕਿ ਇੱਕ ਜੀਵਿਤ ਜੀਵ ਸੁੱਕੀ (ਮੁਰਦਾ) ਲੱਕੜੀ ਵਿਚੋਂ ਸਿਰਜਿਆ ਜਾਂਦਾ ਹੈ। (ਕੋਹਲੀ, ਸੁਰਿੰਦਰ ਸਿੰਘ, ਗੁਰੂ ਗ੍ਰੰਥ ਸਾਹਿਬ ਜੀ ਵਿੱਚ ਮਿਥਿਹਾਸਕ ਹਵਾਲਿਆਂ ਦਾ ਕੋਸ਼, ੧੯੯੩)

%% ---------------- 3 ----------------
\clearpage
\chapter*{\eng{3.} ਅੱਗ ਦਾ ਮੂਲ}
\addcontentsline{toc}{chapter}{\eng{3.} ਅੱਗ ਦਾ ਮੂਲ}
\section*{ਵਿਗਿਆਨਿਕ ਦ੍ਰਿਸ਼ਟੀਕੋਣ}
ਆਮ ਤੌਰ ਤੇ, ਵਾਯੂਮੰਡਲ ਵਿੱਚ ਆਕਸੀਜਨ ਅਤੇ ਕਿਸੇ ਕਿਸਮ ਦੇ ਬਾਲਣ (ਉਦਾਹਰਨ ਲਈ ਲੱਕੜੀ ਜਾਂ ਗੈਸੋਲਿਨ) ਆਪਣੇ ਆਪ ਹੀ ਅੱਗ ਨਹੀਂ ਫੜਦੇ ਕਿਉਂਕਿ ਉਹ ਆਕਸੀਜਨ ਨਾਲ ਘਿਰੇ ਹੁੰਦੇ ਹਨ। ਬਲਣ ਪ੍ਰਤੀਕਿਰਿਆ ਵਾਪਰਨ ਲਈ, ਬਾਲਣ ਨੂੰ ਇਸ ਦੇ ਇਗਨੀਸ਼ਨ ਤਾਪਮਾਨ ਤੱਕ ਗਰਮ ਕਰਨਾ ਜ਼ਰੂਰੀ ਹੈ। ਜਦੋਂ ਲੱਕੜੀ ਲਗਭਗ ੩੦੦\eng{°F} (੧੫੦\eng{°C}) ਤਕ ਪਹੁੰਚਦੀ ਹੈ, ਗਰਮੀ ਕੁਝ ਕੋਮਲ ਮੱਗਰੀ ਨੂੰ ਸਾੜ੍ਹ ਦਿੰਦੀ ਹੈ, ਧੂੰਆਂ (ਹਾਈਡਰੋਜਨ, ਕਾਰਬਨ ਅਤੇ ਆਕਸੀਜਨ ਦਾ ਮਿਸ਼ਰਨ) ਛੱਡਦੀ ਹੈ, ਅਤੇ ਬਾਕੀ ਸਮੱਗਰੀ ਕੋਲਾ ਅਤੇ ਸੁਆਹ ਬਣਾਉਂਦੀ ਹੈ। ਲਾਟ ਦਾ ਰੰਗ ਤਾਪਮਾਨ ਤੇ ਨਿਰਭਰ ਕਰਦਾ ਹੈ — ਸਭ ਤੋਂ ਗਰਮ ਹਿੱਸਾ (ਬੇਸ) ਨੀਲਾ ਚਮਕਦਾ ਹੈ, ਅਤੇ ਉੱਪਰਲੇ ਠੰਡੇ ਹਿੱਸੇ ਸੰਤਰੀ ਜਾਂ ਪੀਲੇ ਚਮਕਦੇ ਹਨ। ਧਰਤੀ ਦੀ ਗੁਰੂਤਾ ਆਕਰਸ਼ਨ ਇਹ ਨਿਰਧਾਰਿਤ ਕਰਦੀ ਹੈ ਕਿ ਲਾਟ ਉੱਪਰ ਵੱਲ ਹੀ ਕਿਉਂ ਉੱਠਦੀ ਹੈ, ਕਿਉਂਕਿ ਗਰਮ ਗੈਸਾਂ ਆਲੇ-ਦੁਆਲੇ ਦੀ ਠੰਢੀ ਹਵਾ ਨਾਲੋਂ ਹਲਕੀਆਂ ਹੁੰਦੀਆਂ ਹਨ।

\section*{ਆਯੁਰਵੈਦਿਕ ਦ੍ਰਿਸ਼ਟੀਕੋਣ — ਤੱਤ: ਜੀਵਨ ਦੇ ਬਿਲਡਿੰਗ ਬਲਾਕ}
ਸ੍ਰਿਸ਼ਟੀ ਦੇ ਸਿਧਾਂਤਾਂ ਨੇ ਪੰਜ ਜ਼ਰੂਰੀ ਤੱਤ ਵਿਕਸਿਤ ਕੀਤੇ ਜੋ ਸਾਰੇ ਜੀਵ ਰੂਪਾਂ ਵਿੱਚ ਸ਼ਾਮਲ ਹਨ: ਈਥਰ, ਹਵਾ, ਅੱਗ, ਪਾਣੀ ਅਤੇ ਧਰਤੀ। ਸਭ ਤੋਂ ਸੂਖਮ ਰੂਪ ਈਥਰ ਹੈ; ਈਥਰ ਅਤੇ ਘਸਰ ਦੇ ਰਲੇਵੇਂ ਤੋਂ ਹਵਾ ਬਣਦੀ ਹੈ; ਹਵਾ ਦੀ ਗਤੀ ਤੋਂ ਰਗੜ੍ਹ ਪੈਦਾ ਹੁੰਦੀ ਹੈ ਜੋ ਗਰਮੀ (ਅੱਗ) ਪੈਦਾ ਕਰਦੀ ਹੈ; ਗਰਮੀ ਨਮੀ ਪੈਦਾ ਕਰਦੀ ਹੈ, ਇਸ ਤਰ੍ਹਾਂ ਪਾਣੀ ਬਣਦਾ ਹੈ; ਅਤੇ ਪਾਣੀ ਦੇ ਸੰਘਣੇਪਣ ਤੋਂ ਧਰਤੀ ਬਣਦੀ ਹੈ। ਵੇਦ ਕਹਿੰਦੇ ਹਨ ਕਿ ਮਨੁੱਖਾਂ ਸਮੇਤ ਸਾਰੀ ਸ੍ਰਿਸ਼ਟੀ ਇਹਨਾਂ ਪੰਜ ਜ਼ਰੂਰੀ ਤੱਤਾਂ ਦੇ ਸੁਮੇਲ ਨਾਲ ਬਣੀ ਹੈ।

\section*{ਮਨੁੱਖੀ ਸਰੀਰ ਵਿੱਚ ਅੱਗ — ਪਿੱਤ}
ਮਨੁੱਖੀ ਸਰੀਰ ਵਿੱਚ ਅੱਗ ਛੇ ਰੂਪਾਂ ਵਿੱਚ ਮੌਜੂਦ ਹੈ, ਜਿਸ ਨੂੰ ਆਮ ਤੌਰ ਤੇ ਪਿੱਤ ਕਿਹਾ ਜਾਂਦਾ ਹੈ:
\begin{enumerate}[label=(\roman*)]
    \item ਪਾਚਕ ਪਿੱਤ — ਛੋਟੀ ਆਂਦਰ, ਪੇਟ ਅਤੇ ਕੋਲਨ ਵਿੱਚ ਪਾਚਨ ਅੱਗ ਵਜੋਂ ਮੌਜੂਦ ਹੈ, ਜੋ ਭੋਜਨ ਨੂੰ ਹਜ਼ਮ ਅਤੇ ਰੂਪਾਂਤਰਿਤ ਕਰਦੀ ਹੈ।
    \item ਪਾਚਕ ਆਪਣੇ ਆਪ ਦੂਜੇ ਚਾਰ ਪਿੱਤਿਆਂ ਨੂੰ ਪੋਸ਼ਣ ਦਿੰਦਾ ਹੈ।
    \item ਰੰਜਕਾ — ਜਿਗਰ ਅਤੇ ਤਿੱਲੀ ਵਿਚ ਸਥਿਤ, ਖੂਨ ਵਿੱਚ ਬਦਲਦਾ ਹੈ।
    \item ਸਾਧਕ — ਹਿਰਦੇ ਵਿੱਚ ਸਥਿਤ, ਗਿਆਨ, ਬੁੱਧੀ ਅਤੇ ਚੇਤਨਾ ਵਿਚ ਮਦਦ ਕਰਦਾ ਹੈ।
    \item ਆਲੋਚਕਾ — ਅੱਖਾਂ ਵਿੱਚ ਸਥਿਤ, ਨਜ਼ਰ ਨੂੰ ਨਿਯੰਤਰਿਤ ਕਰਦਾ ਹੈ।
    \item ਭ੍ਰਾਜਿਕਾ — ਚਮੜੀ ਵਿੱਚ ਵਸਦਾ ਹੈ, ਰੰਗ ਨੂੰ ਨਿਯੰਤਰਿਤ ਕਰਦਾ ਹੈ।
\end{enumerate}
ਪਿੱਤ ਦੀ ਆਪਣੀ ਸ਼ਕਤੀ ਅਕਾਲ ਪੁਰਖ ਵਾਹਿਗੁਰੂ ਜੀ ਤੋਂ ਪ੍ਰਾਪਤ ਹੁੰਦੀ ਹੈ, ਜੋ ਜੋਤਿ ਰੂਪ ਵਿੱਚ ਮਨੁੱਖੀ ਸਰੀਰ ਵਿੱਚ ਨਿਵਾਸ ਕਰਦੇ ਹਨ।

%% ---------------- 4 ----------------
\clearpage
\chapter*{\eng{4.} ਗੁਰਮਤਿ ਦੇ ਚਾਨਣੇ ਵਿੱਚ ਹਵਨ}
\addcontentsline{toc}{chapter}{\eng{4.} ਗੁਰਮਤਿ ਦੇ ਚਾਨਣੇ ਵਿੱਚ ਹਵਨ}
ਗੁਰੂ ਮਰਿਯਾਦਾ ਅਨੁਸਾਰ ਅਸੀਂ ਸਾਰੇ ਭੋਜਨ ਦੇ ਲੰਗਰ ਤੋਂ ਜਾਣੂ ਹਾਂ। ਇਹ ਇੱਕ ਕਿਸਮ ਦਾ ਹਵਨ ਹੈ, ਜਿੱਥੇ ਗੁਰੂ ਖਾਲਸਾ ਪੰਥ ਸਾਧ ਸੰਗਤ (ਪਵਿੱਤਰ ਸੰਗਤ) ਨੂੰ ਭੋਜਨ ਦੇਣ ਲਈ ਅਗਨੀ (ਅੱਗ) ਦੀ ਵਰਤੋਂ ਕਰਦਾ ਹੈ। ਇਥੋਂ ਅੱਗ ਦੀ ਮਹੱਤਤਾ ਨੂੰ ਆਸਾਨੀ ਨਾਲ ਸਥਾਪਿਤ ਕੀਤਾ ਜਾ ਸਕਦਾ ਹੈ, ਕਿਉਂਕਿ ਅਗਨੀ ਤੋਂ ਬਿਨਾ ਭੋਜਨ ਨਹੀਂ ਪਕਾਇਆ ਜਾ ਸਕਦਾ। ਲੰਗਰ ਸੇਵਾ ਦਾ ਇੱਕ ਸਾਲ ੧੨ ਸਾਲਾਂ ਦੇ ਹਵਨ ਦੇ ਬਰਾਬਰ ਹੈ (ਜੋ ਕਿ ਇੱਕ ਕੁੰਡ ਵਿੱਚ ਸਥਾਪਿਤ ਕੀਤਾ ਜਾਂਦਾ ਹੈ, ਅਤੇ ਅੱਗ ਵਿੱਚ ਚੜ੍ਹਾਵਾ ਚੜ੍ਹਾਇਆ ਜਾਂਦਾ ਹੈ)।

ਜ਼ਿਕਰਯੋਗ ਹੈ ਕਿ ਜੋ ਭੋਜਨ ਆਮ ਤੌਰ ਤੇ ਹੋਟਲਾਂ ਅਤੇ ਹੋਰ ਖਾਣ ਪੀਣ ਵਾਲੀਆਂ ਥਾਂਵਾਂ ਤੇ ਪਰੋਸਿਆ ਜਾਂਦਾ ਹੈ, ਉਸਨੂੰ ਲੰਗਰ ਨਹੀਂ ਕਿਹਾ ਜਾਂਦਾ। ਕਾਰਨ ਹੈ ਕਿ ਇਹ ਗੁਰਬਾਣੀ (ਪਰਮਾਤਮਾ/ਗੁਰੂ ਦਾ ਸ਼ਬਦ) ਹੈ ਜੋ ਤਿਆਰ ਕੀਤੇ ਜਾ ਰਹੇ ਭੋਜਨ ਨੂੰ ਪ੍ਰਸ਼ਾਦਿ (ਪਵਿੱਤਰ ਭੋਜਨ) ਵਿੱਚ ਬਦਲ ਦਿੰਦਾ ਹੈ। ਇਸ ਲਈ ਪੜ੍ਹ ਕੇ ਅਸੀਂ ਗੁਰਬਾਣੀ ਦੀ ਭਾਵਨਾ ਨੂੰ ਭੋਜਨ ਵਿੱਚ ਪ੍ਰਫੁੱਲਤ ਕਰਦੇ ਹਾਂ।

ਦੂਸਰੀ ਕਿਸਮ ਦਾ ਹਵਨ, ਜਿਸ ਨੂੰ ਖਾਲਸਾ ਪੰਥ ਨੇ ਸੀਮਿਤ ਅਤੇ ਖੋਜ ਕਾਰਨ ਘਟਦਾ ਦੇਖਿਆ ਹੈ, ਉਹ ਹੈ ਅੱਗ/ਚੰਡਾ ਲਗਾਉਣਾ ਅਤੇ ਉਸ ਦੇ ਦੁਆਲੇ/ਨਾਲ ਬੈਠ ਕੇ ਗੁਰਬਾਣੀ ਪੜ੍ਹਨਾ। ਅੱਜ ਬਹੁਤੇ ਸਿੱਖ ਇਸ ਨੂੰ ਹਿੰਦੂ ਰਸਮ ਕਹਿੰਦੇ ਹਨ। ਪਰ ਇਹ ਕਿਹਾ ਜਾ ਰਿਹਾ ਹੈ ਕਿ ਅਸੀਂ ਗੁਰਬਾਣੀ ਪੜ੍ਹਦੇ ਸਮੇਂ ਬਰਤਨ ਵਿੱਚ ਪਾਣੀ ਰੱਖਦੇ ਹਾਂ, ਅਸੀਂ ਗੁਰੂ ਗ੍ਰੰਥ ਸਾਹਿਬ ਜੀ ਮਹਾਰਾਜ ਨੂੰ ਫੁੱਲ ਅਤੇ ਧੂਪ ਚੜ੍ਹਾਉਂਦੇ ਹਾਂ, ਅਸੀਂ ਚੌਰ ਸਾਹਿਬ ਅਤੇ ਹੋਰ ਬਹੁਤ ਸਾਰੇ ਕਰਮ ਕਾਂਡ ਕਰਦੇ ਹਾਂ। ਰੀਤੀ ਰਿਵਾਜਾਂ ਦਾ ਆਪਣੇ ਆਪ ਵਿੱਚ ਕੋਈ ਅਰਥ ਨਹੀਂ ਹੈ, ਜੇਕਰ ਇਹ ਪ੍ਰੇਮ ਭਾਵਨਾ ਨਾਲ ਨਹੀਂ ਕੀਤੇ ਜਾਂਦੇ। ਭਗਤੀ ਕੇਵਲ ਪਿਆਰ ਤੋਂ ਪੈਦਾ ਹੁੰਦੀ ਹੈ। ਇਸ ਤੋਂ ਇਲਾਵਾ ਸਿੱਖ ਸ੍ਰੀ ਅਕਾਲ ਪੁਰਖ ਸਾਹਿਬ ਜੀ ਦੀ ਪੂਜਾ ਕਰਦੇ ਹਨ, ਨਾ ਕਿ ਦੇਵਤਿਆਂ ਦੀ; ਹਾਲਾਂਕਿ ਅਸੀਂ ਉਹਨਾਂ ਦਾ ਸਤਿਕਾਰ ਕਰਦੇ ਹਾਂ ਕਿਉਂਕਿ ਉਹ ਸ੍ਰੀ ਵਾਹਿਗੁਰੂ ਜੀ ਦੀ ਰਚਨਾ ਹਨ।

\section*{ਸ੍ਰੀ ਦਸਮ ਗ੍ਰੰਥ ਸਾਹਿਬ ਜੀ ਵਿਚੋਂ ਪ੍ਰਮਾਣ}
ਹੇਠਾਂ ਸ੍ਰੀ ਦਸਮ ਗ੍ਰੰਥ ਸਾਹਿਬ ਜੀ ਦਾ ਇੱਕ ਲੇਖ ਹੈ ਜਿੱਥੇ ਗੁਰੂ ਗੋਬਿੰਦ ਸਿੰਘ ਜੀ ਖਾਲਸੇ ਨੂੰ ਉਪਦੇਸ਼ ਦਿੰਦੇ ਹਨ। ਇਹ ਮੰਗਲਾਚਰਨ ਦੀ ਪਾਲਣਾ ਕਰਦਾ ਹੈ ਅਤੇ ਸ੍ਰੀ ਦਸਮ ਗ੍ਰੰਥ ਸਾਹਿਬ ਜੀ ਦੀ ਮੁੱਖ ਕਥਾ ਤੋਂ ਠੀਕ ਪਹਿਲਾਂ ਹੈ।

\begin{quote}
ਪੂਜਾ, ਤਿਲਕ, ਹੋਮ, ਗਾਯਤਰੀ, ਸੰਧਯਾ, ਤਰਪਨ ਧਾਰੋ ॥ \\
\textbf{ਅਰਥ:} ਭਗਤੀ, ਤਿਲਕ, ਹਵਨ, ਗਾਇਤਰੀ ਮੰਤਰ, ਤੀਰਥਾਂ ਤੇ ਭਗਤੀ — ਇਹ ਸਭ ਗੁਣ ਹਿਰਦੇ ਚ ਧਾਰੋ।

ਕੀਟ ਹੋਤਿ ਮੇਰੂ, ਤ੍ਰਿਣ ਮਹਿ ਬਸਿ ਨਿਰੰਜਨਿ ਜਾਨ ॥ \\
ਸਤਯ ਰੂਪ ਆਤਮ ਅਬਿਨਾਸੀ, ਬ੍ਰਹਮ ਤਿਸ ਪਹਿਚਾਨ ॥ \\
ਏਕ ਬ੍ਰਹਮ ਸਭ ਘਟਿ ਘਟਿ ਪੂਰਨ, ਆਦਿ ਪੁਰਖੁ ਭਗਵਾਨ ॥੧੦॥ \\
\textbf{ਅਰਥ:} ਇੱਕ ਛੋਟੀ ਕੀੜੀ ਤੋਂ ਲੈ ਕੇ ਵੱਡੇ ਪਹਾੜ ਅਤੇ ਘਾਹ ਦੇ ਇੱਕ ਛੋਟੇ ਬਲੇਡ ਵਿੱਚ ਵੀ ਪ੍ਰਭੂ ਨਿਰੰਜਨ ਸਭ ਵਿਚ ਵੱਸਦਾ ਹੈ। ਤੁਹਾਡਾ ਸੱਚਾ ਆਤਮ-ਰੂਪ ਕਦੇ ਨਹੀਂ ਮਰਦਾ, ਬ੍ਰਹਮ ਪਰਮਾਤਮਾ ਨੂੰ ਹਮੇਸ਼ਾ ਸੱਚਾ ਮੰਨੋ। ਇੱਕ ਬ੍ਰਹਮ ਸਭ ਵਿੱਚ ਹੈ; ਭਗਵਾਨ ਸ਼ੁਰੂ ਤੋਂ ਹੀ ਇੱਥੇ ਸਨ।
\end{quote}

ਗੁਰਬਾਣੀ ਅਨੁਸਾਰ ਹਵਨ/ਹੋਮ/ਯੱਗ ਦੇ ਹਵਾਲੇ ਸ੍ਰੀ ਗੁਰੂ ਗ੍ਰੰਥ ਸਾਹਿਬ, ਸ੍ਰੀ ਦਸਮ ਗ੍ਰੰਥ ਸਾਹਿਬ, ਸ੍ਰੀ ਸਰਬਲੋਹ ਗ੍ਰੰਥ ਸਾਹਿਬ, ਭਾਈ ਗੁਰਦਾਸ ਜੀ, ਸੂਰਜ ਪ੍ਰਕਾਸ਼, ਪੰਥ ਪ੍ਰਕਾਸ਼ ਅਤੇ ਹੋਰ ਬਹੁਤ ਸਾਰੇ ਸਿੱਖ ਗ੍ਰੰਥਾਂ ਵਿੱਚ ਮਿਲਦੇ ਹਨ। ਕੁਝ ਪ੍ਰਮੁੱਖ ਹਵਾਲੇ:
\begin{itemize}
    \item ਵੁਠੈ ਹੋਇ ਅਨੰਦੁ ... ਤਿਤੁ ਘਿਇ ਹੋਮ ਜਗ ਸਦ ਪੂਜਾ ਪਇਐ ਕਾਰਜੁ ਹੋਇ ॥ (ਆਦਿ ੧੫੦)
    \item ਹੋਮ ਜਗ ਅਰੁ ਪਾਠ ਪੁਰਾਣ ॥ ਜੋ ਤਿਸੁ ਭਾਵੈ ਸੋ ਪਰਵਾਨੁ ॥੩॥ (ਆਦਿ ੧੨੫੭)
    \item ਭਾਉ ਭਗਤਿ ਪਰਵਾਣ ਹੈ, ਜਗ ਹੋਮ ਤੇ ਪੁਰਬ ਕਮਾਈ ॥ (ਭਾਈ ਗੁਰਦਾਸ ਜੀ)
    \item ਪੂਜਾ ਤਿਲਕ ਹੋਮ ਗਾਯਤ੍ਰੀ ਸੰਧਯਾ ਤਰਪਨ ਧਾਰੋ ॥ (ਸ੍ਰੀ ਸਰਬਲੋਹ ਗ੍ਰੰਥ ਸਾਹਿਬ)
\end{itemize}

ਸ਼ਹੀਦ ਬਾਬਾ ਦੀਪ ਸਿੰਘ ਜੀ ਦੇ ਜੰਗ ਲਈ ਜਾਣ ਤੋਂ ਪਹਿਲਾਂ ਦੀ ਘਟਨਾ ਵਿੱਚ ਵੀ ਹਵਨ ਦਾ ਵਰਣਨ ਮਿਲਦਾ ਹੈ — ਬਾਬਾ ਜੀ ਨੇ ਇੱਕ ਅਖੰਡ ਪਾਠ ਦਾ ਆਯੋਜਨ ਕੀਤਾ ਅਤੇ ਇੱਕ ਹਵਨ ਪੂਰਾ ਕੀਤਾ (ਲੜਾਈ ਵੱਲ ਜਾਣ ਤੋਂ ਪਹਿਲਾਂ) ਅਤੇ ਆਪਣੇ ਗੁੱਟ ਦੇ ਦੁਆਲੇ ਇੱਕ ਵਿਆਹ ਦਾ ਕੰਗਣ ਬੰਨ੍ਹਿਆ (ਸ਼ਹਾਦਤ ਦੀ ਤਿਆਰੀ ਵਜੋਂ ਮੌਤ ਨੂੰ ਪ੍ਰਭੂ ਦੇ ਨਾਲ ਅਭੇਦ ਹੋਣਾ ਦੇਖਿਆ ਜਾਂਦਾ ਹੈ)।

ਇਸ ਕਥਨ ਅਨੁਸਾਰ ਹਵਨ ਦੇ ਦੋਵੇਂ ਅਰਥ ਲਾਗੂ ਕੀਤੇ ਜਾ ਸਕਦੇ ਹਨ ਅਤੇ ਦੋਵੇਂ ਹੀ ਸਹੀ ਜਾਪਦੇ ਹਨ ਕਿਉਂਕਿ ਯੁੱਧ ਲਈ ਆਪਣੇ ਆਪ ਨੂੰ ਤਿਆਰ ਕਰਨ ਲਈ ਸਰੀਰਕ ਅਤੇ ਮਾਨਸਿਕ ਤੌਰ 'ਤੇ ਤੰਦਰੁਸਤ ਹੋਣਾ ਜ਼ਰੂਰੀ ਹੈ।

%% ---------------- 5 ----------------
\clearpage
\chapter*{\eng{5.} ਹਵਨ ਦੇ ਵਿਗਿਆਨਕ ਲਾਭ}
\addcontentsline{toc}{chapter}{\eng{5.} ਹਵਨ ਦੇ ਵਿਗਿਆਨਕ ਲਾਭ}
ਹਵਨ (ਹੋਮ) ਇੱਕ ਰਸਮ ਨੂੰ ਦਰਸਾਉਂਦਾ ਹੈ, ਜਿਸ ਵਿੱਚ ਇੱਕ ਬਲੀ ਜਾਂ ਕੋਈ ਧਾਰਮਿਕ ਭੇਟ ਅੱਗ ਵਿੱਚ ਕੀਤੀ ਜਾਂਦੀ ਹੈ। ਅੱਗ ਏਜੰਟ ਹੈ, ਅਤੇ ਭੇਟਾਂ ਵਿੱਚ ਉਹ ਸ਼ਾਮਲ ਹਨ ਜੋ ਪਦਾਰਥਕ ਅਤੇ ਪ੍ਰਤੀਕ ਹਨ ਜਿਵੇਂ ਕਿ ਅਨਾਜ, ਘਿਓ, ਦੁੱਧ, ਧੂਪ ਅਤੇ ਬੀਜ।

ਜਦੋਂ ਹਵਨ ਕੀਤਾ ਜਾਂਦਾ ਹੈ ਤਾਂ ਹਵਾ ਅੱਗ ਨੂੰ ਛੂੰਹਦੀ ਹੈ, ਸ਼ੁੱਧ ਹੋ ਜਾਂਦੀ ਹੈ ਅਤੇ ਤਾਸੀਰ ਵਿੱਚ ਹਲਕੀ ਹੋ ਜਾਂਦੀ ਹੈ, ਉੱਪਰ ਜਾਂਦੀ ਹੈ ਅਤੇ ਠੰਢੀ/ਅਸ਼ੁੱਧ ਹਵਾ ਦੀ ਥਾਂ ਲੈਂਦੀ ਹੈ। ਹਵਨ ਵਿੱਚ ਮੁੱਖ ਸਮੱਗਰੀ ਅੰਬ ਦੀ ਲੱਕੜੀ ਹੁੰਦੀ ਹੈ ਜਿਸ ਨੂੰ ਸਾੜ੍ਹਨ 'ਤੇ ਫਾਰਮਿਕ ਐਲਡੀਹਾਈਡ ਗੈਸ ਨਿਕਲਦੀ ਹੈ ਜੋ ਹਾਨੀਕਾਰਕ ਬੈਕਟੀਰੀਆ ਨੂੰ ਮਾਰ ਦਿੰਦੀ ਹੈ, ਇਸ ਤਰ੍ਹਾਂ ਵਾਤਾਵਰਣ ਨੂੰ ਸ਼ੁੱਧ ਕਰਦਾ ਹੈ। ਗਾਂ ਦਾ ਘਿਓ ਹਵਨ ਲਈ ਬਹੁਤ ਮਹੱਤਵਪੂਰਨ ਸਮੱਗਰੀ ਹੈ ਅਤੇ ਇਸ ਦੀ ਮਹਿਕ ਵਾਤਾਵਰਨ ਨੂੰ ਸ਼ੁੱਧ ਕਰਦੀ ਹੈ। ਇਹ ਯਗਿਆ ਦੀ ਵਿਗਿਆਨਕ ਵਿਆਖਿਆ ਦੇ ਰੂਪ ਵਿੱਚ ਪ੍ਰਕਾਸ਼ਿਤ ਕੀਤਾ ਗਿਆ ਸੀ (\eng{26-28} ਜੁਲਾਈ \eng{1996} ਨੂੰ ਮਾਂਟਰੀਅਲ, ਕੈਨੇਡਾ ਵਿੱਚ ਆਯੋਜਿਤ ਅਸ਼ਵਮੇਧ ਯਗਿਆ ਦੀ ਕਾਰਵਾਈ ਵਿੱਚ ਪ੍ਰਕਾਸ਼ਿਤ ਇੱਕ ਲੇਖ)।

ਹਵਨ ਦੀ ਪ੍ਰਕਿਰਿਆ ਦੌਰਾਨ ਹਵਨ ਸਮੱਗਰੀ ਨੂੰ ਸਾੜ੍ਹਨ ਵੇਲੇ ਜੋ ਰਸਾਇਣਕ ਕਿਰਿਆ ਹੁੰਦੀ ਹੈ ਉਹ ਹੈ:
\begin{center}
\eng{CO\textsubscript{2} + H\textsubscript{2}O + 112,000 cal $\rightarrow$ HCHO + O\textsubscript{2}}
\end{center}
(ਕਾਰਬਨ ਡਾਈਆਕਸਾਈਡ + ਪਾਣੀ + ਊਰਜਾ $\rightarrow$ ਫਾਰਮਾਲਡੀਹਾਈਡ + ਆਕਸੀਜਨ)। ਇਸ ਲਈ ਹਵਨ ਦੁਆਰਾ ਕਾਫੀ ਹੱਦ ਤੱਕ ਕਾਰਬਨ ਡਾਈਆਕਸਾਈਡ ਨੂੰ ਆਕਸੀਜਨ ਵਿੱਚ ਬਦਲ ਦਿੱਤਾ ਜਾਂਦਾ ਹੈ।

ਡਾ. ਕੁੰਦਨ ਲਾਲ ਐਮ.ਡੀ ਦੁਆਰਾ ਕੀਤੀ ਖੋਜ ਅਨੁਸਾਰ ਇਹ ਪਾਇਆ ਗਿਆ ਕਿ ਜਦੋਂ ਹਵਨ ਮਾਧਿਅਮ ਵਿੱਚ ਅੰਬ ਦੀ ਲੱਕੜੀ ਸਾੜ੍ਹੀ ਗਈ ਤਾਂ ਬੈਕਟੀਰੀਆ ਦੀ ਗਿਣਤੀ \eng{94\%} ਤੱਕ ਘੱਟ ਗਈ, ਅਤੇ \eng{24} ਘੰਟਿਆਂ ਬਾਅਦ ਵੀ \eng{96\%} ਘੱਟ ਰਹੀ। ਹਵਨ ਦਾ ਧੂੰਆਂ ਆਯੁਰਵੇਦ ਤੋਂ ਲਿਆ ਗਿਆ ਇਨਹੇਲੇਸ਼ਨ ਥੈਰੇਪੀ ਵਰਗਾ ਹੈ ਅਤੇ ਮਨ ਨੂੰ ਅਰਾਮ ਦਿੰਦਾ ਹੈ ਅਤੇ ਸਰੀਰ ਨੂੰ ਚੰਗਾ ਕਰਦਾ ਹੈ; ਜਦੋਂ ਹਵਨ ਨੂੰ ਗੁਰਬਾਣੀ ਦੇ ਪਾਠ ਨਾਲ ਜੋੜਿਆ ਜਾਂਦਾ ਹੈ, ਤਾਂ ਹਵਨ ਕਰਨ ਵਾਲੇ ਅਤੇ ਹਵਨ ਕਰਵਾਉਣ ਵਾਲੇ ਇੱਕ ਡੂੰਘੀ ਅਤੇ ਲੰਮੀ ਸਮਾਧੀ ਦਾ ਅਨੁਭਵ ਕਰਦੇ ਹਨ।

ਜੇਕਰ ਕੋਈ ਵੀ ਰਸਮ (ਨੇਮ) ਗੁਰਬਾਣੀ ਅਭਿਆਸ ਦੇ ਉਦੇਸ਼, ਵਿਸ਼ਵਾਸ ਅਤੇ ਪਿਆਰ, ਅਤੇ ਗਿਆਨ — ਇਹਨਾਂ ਤਿੰਨਾਂ ਬਿੰਦੂਆਂ ਵਿੱਚੋਂ ਕਿਸੇ ਤੋਂ ਵੀ ਸੱਖਣੀ ਹੋਵੇ ਤਾਂ ਉਸ ਦਾ ਪੂਰਾ ਲਾਭ ਨਹੀਂ ਮਿਲੇਗਾ।

%% ---------------- 6 ----------------
\clearpage
\chapter*{\eng{6.} ਹਵਨ ਦਾ ਭੌਤਿਕ ਪੱਖ}
\addcontentsline{toc}{chapter}{\eng{6.} ਹਵਨ ਦਾ ਭੌਤਿਕ ਪੱਖ}
ਯੱਗ ਦੇ ਮੂਲ ਬਲੀਦਾਨ ਪ੍ਰਦਰਸ਼ਨ ਵਿੱਚ ਛੁਪੇ ਆਦਰਸ਼ਾਂ 'ਤੇ ਵਿਚਾਰ ਕਰੀਏ ਤਾਂ ਸਭ ਦਾ ਕਲਿਆਣ ਇਸ ਦੇ ਮੁੱਖ ਉਦੇਸ਼ ਵਜੋਂ ਪ੍ਰਗਟ ਹੁੰਦਾ ਹੈ। ਜਿਵੇਂ ਕਿ ਭਗਵਾਨ ਕ੍ਰਿਸ਼ਨ ਦੁਆਰਾ ਗੀਤਾ ਵਿੱਚ ਜ਼ਿਕਰ ਕੀਤਾ ਗਿਆ ਹੈ (ਅਧਿਆਇ ੩, ਸ਼ਲੋਕ ੯-੧੫): ਯੱਗ ਹਵਾ ਨੂੰ ਸ਼ੁੱਧ ਕਰਦਾ ਹੈ, ਜਿਸ ਨਾਲ ਬਾਰਿਸ਼ ਹੁੰਦੀ ਹੈ, ਸਾਰੀ ਬਨਸਪਤੀ ਅਤੇ ਜੜ੍ਹੀ-ਬੂਟੀਆਂ ਨੂੰ ਸ਼ੁੱਧ ਕਰਦਾ ਹੈ, ਅਤੇ ਅੰਤ ਵਿੱਚ ਅਨਾਜ ਪੈਦਾ ਕਰਦਾ ਹੈ ਜੋ ਜੀਵਨ ਦੀ ਲੰਬੀ ਉਮਰ ਪ੍ਰਦਾਨ ਕਰਦਾ ਹੈ।

ਹਵਨ ਇੱਕ ਪ੍ਰਾਚੀਨ ਰਸਮ ਹੈ ਜੋ ਵਾਤਾਵਰਣ ਨੂੰ ਸ਼ੁੱਧ ਕਰਨ ਲਈ ਕੀਤੀ ਜਾਂਦੀ ਹੈ। ਕਿਉਂਕਿ ਫਾਰਮਲਡੀਹਾਈਡ ਸਿਰਫ ਪਾਣੀ ਦੀ ਮੌਜੂਦਗੀ ਵਿੱਚ ਪ੍ਰਭਾਵਸ਼ਾਲੀ ਹੁੰਦਾ ਹੈ, ਇਸ ਲਈ ਹਵਨ ਕੁੰਡ ਦੇ ਆਲੇ ਦੁਆਲੇ ਅਤੇ ਹਵਾ ਵਿੱਚ ਪਾਣੀ ਛਿੜਕਣ ਦੀ ਰਸਮ ਹੈ। “ਹਵਨ ਨਾਲ ਮਾਹੌਲ ਨੂੰ ਠੀਕ ਕਰੋ ਅਤੇ ਚੰਗਾ ਮਾਹੌਲ ਤੁਹਾਨੂੰ ਠੀਕ ਕਰ ਦੇਵੇਗਾ” — ਡਾ. ਮਧੂਕਰ ਗਿਆਕਵਾੜ। ਡਾ. ਹਾਫਕਿਨ ਦਾ ਕਹਿਣਾ ਹੈ ਕਿ ਦੂਰੋਂ ਹਵਨ ਦੇ ਧੂੰਏਂ ਨੂੰ ਸਾਹ ਲੈਣ ਨਾਲ ਵਿੰਡਪਾਈਪ ਨਾਲ ਸੰਬੰਧਤ ਕੁਝ ਗ੍ਰੰਥੀਆਂ ਤੋਂ ਸਿਕਰੇਸ਼ਨ ਹੁੰਦਾ ਹੈ ਜੋ ਸਾਡੇ ਮਨ ਅਤੇ ਦਿਲ ਨੂੰ ਆਰਾਮ ਨਾਲ ਭਰ ਦਿੰਦਾ ਹੈ।

ਹਵਨ ਦੀ ਰਸਮ ਕਰਦੇ ਸਮੇਂ ਇਸ ਗੱਲ ਦਾ ਧਿਆਨ ਰੱਖਣਾ ਚਾਹੀਦਾ ਹੈ ਕਿ ਹਵਨ ਦਾ ਵੱਧ ਤੋਂ ਵੱਧ ਲਾਭ ਲੈਣ ਲਈ ਹਵਨ ਸਮੱਗਰੀ ਦਾ ਪੂਰਨ ਦਹਿਨ (ਪੂਰੀ ਤਰ੍ਹਾਂ ਸੜਨਾ) ਹੋਣਾ ਚਾਹੀਦਾ ਹੈ।

%% ---------------- 7 ----------------
\clearpage
\chapter*{\eng{7.} ਸਨਾਤਨ ਪਰੰਪਰਾ ਅਤੇ ਸਿੱਖ ਵਿਸ਼ਵਾਸ ਅਨੁਸਾਰ ਅਗਨੀ ਦੇਵ}
\addcontentsline{toc}{chapter}{\eng{7.} ਸਨਾਤਨ ਪਰੰਪਰਾ ਅਤੇ ਸਿੱਖ ਵਿਸ਼ਵਾਸ ਅਨੁਸਾਰ ਅਗਨੀ ਦੇਵ}
ਸਨਾਤਨ ਗ੍ਰੰਥਾਂ ਵਿੱਚ ਇਹ ਪ੍ਰਸਿੱਧ ਮੰਨਿਆ ਜਾਂਦਾ ਹੈ ਕਿ ਅਗਨੀ ਦੇਵ (ਇੱਕ ਅਰਧ-ਦੇਵਤਾ) ਅੱਗ ਦੇ ਰੂਪ ਵਿੱਚ ਇਸ਼ਟ ਦੇਵਤਾ ਨੂੰ ਖੁਸ਼ ਕਰਨ ਲਈ ਵਰਤਿਆ ਜਾਂਦਾ ਹੈ ਅਤੇ ਭੇਟ ਕੀਤੀ ਸਮੱਗਰੀ ਨੂੰ ਦੇਵਤੇ ਤੱਕ ਪਹੁੰਚਾਉਂਦਾ ਹੈ। ਇਹ ਸਾਡੇ ਸਨਾਤਨੀ ਭਰਾਵਾਂ ਲਈ ਸੱਚ ਹੈ, ਪਰ ਇਹ ਮੁੱਖ ਕਾਰਨ ਨਹੀਂ ਹੈ ਕਿ ਸਿੱਖ ਹਵਨ ਕਿਉਂ ਕਰਦੇ ਹਨ।

ਸ੍ਰੀ ਗੁਰੂ ਗ੍ਰੰਥ ਸਾਹਿਬ ਜੀ ਨੂੰ ਮੰਨਣ ਵਾਲਿਆਂ ਲਈ, ਅਸੀਂ ਕਿਸੇ ਵੀ ਬਖਸ਼ਿਸ਼ ਲਈ ਦੇਵਤਿਆਂ ਕੋਲ ਨਹੀਂ ਜਾਂਦੇ। ਅਸੀਂ ਕੇਵਲ ਸ੍ਰੀ ਅਕਾਲ ਪੁਰਖ ਵਾਹਿਗੁਰੂ ਜੀ ਨੂੰ ਹੀ ਨਤਮਸਤਕ ਹੁੰਦੇ ਹਾਂ, ਜੋ ਇਸ ਸਭ ਦਾ ਸਿਰਜਣਹਾਰ ਹੈ, ਦੇਵਤਿਆਂ ਦਾ ਵੀ ਸਿਰਜਣਹਾਰ ਹੈ। ਗੁਰਬਾਣੀ ਅਨੁਸਾਰ ਸਾਰੇ ਦੇਵੀ ਦੇਵਤੇ ਮਾਇਆ ਦੀ ਰਚਨਾ ਦੇ ਅੰਦਰ ਹਨ, ਇਸ ਲਈ ਦੇਵਤਿਆਂ ਨੂੰ ਖੁਸ਼ ਕਰਨ ਦੀ ਕੋਸ਼ਿਸ਼ ਸਿੱਖਾਂ ਲਈ ਵਿਅਰਥ ਹੈ; ਸਾਡਾ ਨਿਸ਼ਾਨਾ ਉਹ ਹੈ ਜੋ ਮਾਇਆ ਤੋਂ ਉੱਪਰ ਹੈ — ਸਰਵਉੱਚ ਸਿਰਜਣਹਾਰ।

ਇੱਕ ਅੰਦਰੂਨੀ ਹਵਨ ਵੀ ਹੁੰਦਾ ਹੈ, ਜੋ ਹਰ ਵਾਰ ਜਦੋਂ ਅਸੀਂ ਖਾਂਦੇ ਹਾਂ, ਇਸ ਪ੍ਰਕਿਰਿਆ ਨੂੰ ਪਾਚਨ ਕਿਹਾ ਜਾਂਦਾ ਹੈ। ਭੋਜਨ ਪੇਟ ਵਿੱਚ ਦਾਖਲ ਹੁੰਦਾ ਹੈ ਤਾਂ ਅਨੇਕਾਂ ਐਸਿਡਾਂ ਦੇ ਸੁਮੇਲ ਨਾਲ ਟੁੱਟ ਜਾਂਦਾ ਹੈ, ਜਿਸ ਨਾਲ ਪੇਟ ਦਾ ਅੰਦਰੂਨੀ ਤਾਪਮਾਨ ਵਧ ਜਾਂਦਾ ਹੈ — ਇਸ ਨੂੰ ਗੁਰਬਾਣੀ ਵਿੱਚ ਜਠਰ ਅਗਨੀ ਕਿਹਾ ਗਿਆ ਹੈ। ਇਹ ਅਖੌਤੀ ਅਗਨੀ ਦੇਵਤਾ ਪੇਟ ਵਿੱਚ ਨਿਵਾਸ ਕਰਦਾ ਹੈ, ਭੋਜਨ ਨੂੰ ਸੂਖਮ ਰੂਪ ਵਿੱਚ ਬਦਲਦਾ ਹੈ, ਜੋ ਫਿਰ ਖੂਨ ਦੁਆਰਾ ਸਰੀਰ ਵਿੱਚ ਸੰਚਾਰਿਤ ਹੁੰਦਾ ਹੈ, ਅਤੇ ਸਰੀਰ ਵਿੱਚ ਮੌਜੂਦ ਸਾਰੇ ਦੇਵਤਿਆਂ ਨੂੰ ਪੋਸ਼ਣ ਦਿੰਦਾ ਹੈ।

%% ---------------- 8 ----------------
\clearpage
\chapter*{\eng{8.} ਕੌਣ ਅਤੇ ਕਿਵੇਂ?}
\addcontentsline{toc}{chapter}{\eng{8.} ਕੌਣ ਅਤੇ ਕਿਵੇਂ?}
ਕੋਈ ਵੀ ਵਿਅਕਤੀ ਜੋ ਗੁਰਬਾਣੀ ਸਾਧਨਾ/ਅਭਿਆਸ ਕਰਨਾ ਚਾਹੁੰਦਾ ਹੈ, ਆਪਣੇ ਆਲੇ-ਦੁਆਲੇ ਦੀ ਵਾਈਬਰੇਸ਼ਨਲ ਊਰਜਾ ਨੂੰ ਸੁਧਾਰਨਾ ਚਾਹੁੰਦਾ ਹੈ, ਬੀਮਾਰੀ ਨੂੰ ਠੀਕ ਕਰਨਾ ਚਾਹੁੰਦਾ ਹੈ, ਮਾਹੌਲ ਨੂੰ ਪੋਸ਼ਣ ਦੇਣਾ ਚਾਹੁੰਦਾ ਹੈ — ਮਰਦ ਹੋਵੇ ਜਾਂ ਔਰਤ, ਨੌਜਵਾਨ ਹੋਵੇ ਜਾਂ ਬਜ਼ੁਰਗ, ਪੜ੍ਹਿਆ-ਲਿਖਿਆ ਹੋਵੇ ਜਾਂ ਅਨਪੜ੍ਹ, ਅਮੀਰ ਹੋਵੇ ਜਾਂ ਗ਼ਰੀਬ — ਹਵਨ ਕਰ ਸਕਦਾ ਹੈ। ਲੋੜ ਸਿਰਫ਼ ਇਹ ਹੈ ਕਿ ਇਸ ਨੂੰ ਸਹੀ ਢੰਗ ਨਾਲ ਅਤੇ ਦਿਲੀ ਇਰਾਦੇ ਨਾਲ ਕੀਤਾ ਜਾਵੇ। ਇਹ ਇੱਕ ਊਰਜਾ ਸੰਸਕਾਰ ਹੈ ਅਤੇ ਇਸ ਦੇ ਲਾਭਾਂ ਨੂੰ ਪ੍ਰਾਪਤ ਕਰਨ ਲਈ, ਸਾਧਕ ਨੂੰ ਝੂਠ, ਭ੍ਰਿਸ਼ਟਾਚਾਰ ਅਤੇ ਬੁਰੀਆਂ ਆਦਤਾਂ ਤੋਂ ਬਚਣਾ ਚਾਹੀਦਾ ਹੈ ਅਤੇ ਚੰਗੇ ਨੈਤਿਕ ਮੁੱਲ ਧਾਰਨ ਕਰਨੇ ਚਾਹੀਦੇ ਹਨ।

ਹਵਨ ਕੁੰਡ ਦਾ ਸਭ ਤੋਂ ਵੱਧ ਵਰਤਿਆ ਜਾਣ ਵਾਲਾ ਆਕਾਰ ਚੌਰਸ (ਵਰਗ) ਹੈ, ਜੋ ਤੁਹਾਡੇ ਕੋਲ ਉਪਲਬਧ ਖੇਤਰ ਅਨੁਸਾਰ ਵੱਡਾ ਜਾਂ ਛੋਟਾ ਹੋ ਸਕਦਾ ਹੈ। ਜੇਕਰ ਵੱਡਾ ਹਵਨ ਕੁੰਡ ਸਥਾਪਿਤ ਕਰਨ ਲਈ ਜਗ੍ਹਾ ਨਹੀਂ ਹੈ, ਤਾਂ ਪੋਰਟੇਬਲ ਹਵਨ ਕੁੰਡਾਂ (ਤਰਜੀਹੀ ਤੌਰ 'ਤੇ ਤਾਂਬੇ ਦੇ, ਕਿਉਂਕਿ ਤਾਂਬਾ ਗਰਮੀ ਦਾ ਵਧੀਆ ਸੰਚਾਲਕ ਹੈ) ਦੀ ਵਰਤੋਂ ਕਰਕੇ ਵੀ ਲਾਭ ਪ੍ਰਾਪਤ ਕੀਤਾ ਜਾ ਸਕਦਾ ਹੈ। ਦਲ ਪੰਥਾਂ (ਨਿਹੰਗ ਦਲ) ਵਿੱਚ ਇਸ ਨੂੰ ਪੰਜ-ਨਿਸ਼ਾਨਚੀ ਸਿੰਘਾਂ ਅਤੇ ਪੰਜ ਪਿਆਰਿਆਂ ਸਾਹਿਬਾਨ ਦੀ ਧੂਪ-ਬੱਤੀ ਕਰਨ ਲਈ “ਚਕਰਵਰਤੀ ਹਵਨ” ਜਾਂ “ਚੰਦਾ ਪ੍ਰਗਟ ਕਰਨਾ” ਕਿਹਾ ਜਾਂਦਾ ਹੈ, ਅਤੇ ਇਹ ਸ੍ਰੀ ਰਹਿਰਾਸ ਸਾਹਿਬ ਦੇ ਪਾਠ ਤੋਂ ਬਾਅਦ ਅਤੇ ਸ੍ਰੀ ਆਸਾ ਦੀ ਵਾਰ ਤੋਂ ਪਹਿਲਾਂ ਅੰਮ੍ਰਿਤਵੇਲਾ ਦੌਰਾਨ ਆਰਤੀ ਸਮਾਰੋਹ ਵਿੱਚ ਵਰਤਿਆ ਜਾਂਦਾ ਹੈ।

\section*{ਵੱਡੇ ਹਵਨ (ਚੰਡਾ) ਲਈ ਸਮੱਗਰੀ ਦੀ ਸੂਚੀ}
{\small
\setlength{\tabcolsep}{3pt}
\renewcommand{\arraystretch}{1.3}
\rowcolors{2}{nanak_niwas_soft_peach!35}{white}
\begin{longtable}{p{0.8cm}p{3.9cm}p{2.6cm}p{2.6cm}}
\hline
\rowcolor{nanak_niwas_midnight_violet}
\multicolumn{1}{c}{\color{nanak_niwas_metalic_gold}\textbf{ਸ. ਨੰ}} &
\multicolumn{1}{c}{\color{nanak_niwas_metalic_gold}\textbf{ਸਮੱਗਰੀ}} &
\multicolumn{1}{c}{\color{nanak_niwas_metalic_gold}\textbf{ਮਾਤਰਾ (\eng{1ft$\times$1ft})}} &
\multicolumn{1}{c}{\color{nanak_niwas_metalic_gold}\textbf{ਮਾਤਰਾ (\eng{0.5ft$\times$0.5ft})}} \\
\hline
\endfirsthead
\hline
\rowcolor{nanak_niwas_midnight_violet}
\multicolumn{1}{c}{\color{nanak_niwas_metalic_gold}\textbf{ਸ. ਨੰ}} &
\multicolumn{1}{c}{\color{nanak_niwas_metalic_gold}\textbf{ਸਮੱਗਰੀ}} &
\multicolumn{1}{c}{\color{nanak_niwas_metalic_gold}\textbf{ਮਾਤਰਾ (\eng{1ft$\times$1ft})}} &
\multicolumn{1}{c}{\color{nanak_niwas_metalic_gold}\textbf{ਮਾਤਰਾ (\eng{0.5ft$\times$0.5ft})}} \\
\hline
\endhead
\eng{1} & ਲਾਲ ਜਾਂ ਚਿੱਟੀ ਚੰਦਨ ਦੀ ਲੱਕੜੀ ਜਾਂ ਅੰਬ ਦੀਆਂ ਡੰਡੀਆਂ & \centering \eng{5} ਕਿਲੋਗ੍ਰਾਮ & \centering\arraybackslash \eng{2.5} ਕਿਲੋਗ੍ਰਾਮ \\
\eng{2} & ਛੜ & & \\
\eng{3} & ਮੁਥਰਾ & & \\
\eng{4} & ਭੋਜ ਪੱਤਰ (\eng{betula utilis}) & \centering \eng{10-20} ਪੱਤੇ & \centering\arraybackslash \eng{5-8} ਪੱਤੇ \\
\eng{5} & ਜੌਂ (\eng{barley}) & \centering \eng{1} ਕੱਪ & \centering\arraybackslash \eng{1/2} ਕੱਪ \\
\eng{6} & ਇੰਦਰ ਜੌਂ (\eng{holarrhena antidysentrica}) & \centering \eng{1} ਕੱਪ & \centering\arraybackslash \eng{1/2} ਕੱਪ \\
\eng{7} & ਗੁੱਗਲ (\eng{commiphora wightii}) & \centering \eng{300} ਗਰਾਮ & \centering\arraybackslash \eng{100} ਗਰਾਮ \\
\eng{8} & ਸੁਪਾਰੀ (\eng{areca nut}) & \centering \eng{300} ਗਰਾਮ & \centering\arraybackslash \eng{100} ਗਰਾਮ \\
\eng{9} & ਸੰਧੂਰ (\eng{vermillion}) & \centering \eng{200} ਗਰਾਮ & \centering\arraybackslash \eng{50} ਗਰਾਮ \\
\eng{10} & ਛੁਹਾਰੇ (\eng{dry dates}) & \centering \eng{10} & \centering\arraybackslash \eng{3-5} \\
\eng{11} & ਬਦਾਮ (\eng{almonds}) & \centering \eng{20} & \centering\arraybackslash \eng{5-8} \\
\eng{12} & ਕਿਸ਼ਮਿਸ਼ (\eng{raisins}) & \centering \eng{250} ਗਰਾਮ & \centering\arraybackslash \eng{25-50} ਗਰਾਮ \\
\eng{13} & ਖੋਪਾ (\eng{coconut}) & \centering \eng{1} ਪੂਰਾ & \centering\arraybackslash \eng{1/2} \\
\eng{14} & ਇਲਾਚੀ ਹਰੀ (\eng{cardamom}) & \centering \eng{150} ਗਰਾਮ & \centering\arraybackslash \eng{20-40} ਗਰਾਮ \\
\eng{15} & ਲੌਂਗ (\eng{clove}) & \centering \eng{50} ਗਰਾਮ & \centering\arraybackslash \eng{10} ਗਰਾਮ \\
\eng{16} & ਮੋਸ਼ਕ ਕਪੂਰ (\eng{camphor}) & \centering ਲੋੜ ਅਨੁਸਾਰ & \centering\arraybackslash ਲੋੜ ਅਨੁਸਾਰ \\
\eng{17} & ਤਿਲ (\eng{sesame}) & \centering \eng{250} ਗਰਾਮ & \centering\arraybackslash \eng{50-80} ਗਰਾਮ \\
\eng{18} & ਕੇਸਰ (\eng{saffron}) & \centering \eng{25} ਗਰਾਮ & \centering\arraybackslash \eng{5-10} ਗਰਾਮ \\
\eng{19} & ਖੰਡ (\eng{sugar}) & \centering \eng{250} ਗਰਾਮ & \centering\arraybackslash \eng{50} ਗਰਾਮ \\
\eng{20} & ਘਿਓ (\eng{clarified butter}) & \centering \eng{1-2} ਕਿਲੋਗ੍ਰਾਮ & \centering\arraybackslash \eng{1/4-1/2} ਕਿਲੋਗ੍ਰਾਮ \\
\end{longtable}
}

\section*{ਹਵਨ ਕਰਵਾਉਣ ਦਾ ਸਮਾਂ}
ਇਹ ਅੰਮ੍ਰਿਤਵੇਲਾ (ਸੂਰਜ ਚੜ੍ਹਨ ਤੋਂ ਪਹਿਲਾਂ) ਜਾਂ ਸੂਰਜ ਡੁੱਬਣ ਵੇਲੇ ਕੀਤਾ ਜਾ ਸਕਦਾ ਹੈ। ਚਕਰਵਰਤੀ ਹਵਨ ਲਈ ਕੋਈ ਸਖਤ ਸਮਾਂ ਸੀਮਾ ਨਹੀਂ ਹੈ, ਪਰ ਅੰਮ੍ਰਿਤਵੇਲਾ ਅਤੇ ਸੂਰਜ ਡੁੱਬਣ ਵੇਲੇ ਦੀ ਪਾਲਣਾ ਕਰਨਾ ਭਲੀ ਹੈ।

\section*{ਹਵਨ ਕਿਵੇਂ ਕਰੀਏ}
\begin{enumerate}[label=\arabic*.]
    \item ਹਵਨ ਕੁੰਡ ਅਤੇ ਆਲੇ-ਦੁਆਲੇ ਦੀ ਜਗ੍ਹਾ ਨੂੰ ਸਾਫ਼ ਕਰੋ ਅਤੇ ਹਵਾਦਾਰੀ ਯਕੀਨੀ ਬਣਾਓ।
    \item ਸਾਰੀ ਸਮੱਗਰੀ ਤਿਆਰ ਰੱਖੋ। ਪਾਣੀ ਦੇ ਦੋ ਬਰਤਨ (ਕੁੰਭ) ਰੱਖੋ — ਇੱਕ ਸਾਵਧਾਨੀ ਲਈ, ਦੂਜਾ ਢੱਕਣ ਨਾਲ ਢੱਕਿਆ ਹੋਇਆ ਜੋ ਗੁਰਬਾਣੀ ਦੀਆਂ ਧੁਨਾਂ ਨਾਲ ਅੰਮ੍ਰਿਤ ਰੂਪ ਵਿਚ ਬਦਲ ਜਾਵੇਗਾ।
    \item ਕੁਝ ਮਿੰਟਾਂ ਲਈ ਚੁੱਪ ਦਾ ਅਭਿਆਸ ਕਰੋ ਅਤੇ ਇੱਕ ਸਾਧਾਰਨ ਅਰਦਾਸ ਨਾਲ ਸ਼ੁਰੂ ਕਰੋ।
    \item ਕਪੂਰ ਰੱਖੋ, ਹਵਨ ਕੁੰਡ ਵਿਚ ਅੰਬ ਦੀਆਂ ਡੰਡੀਆਂ ਰੱਖੋ ਤਾਂ ਜੋ ਕੁੰਡ ਵਿੱਚ ਹਵਾਦਾਰੀ ਬਣੀ ਰਹੇ।
    \item ਮਾਚਿਸ ਦੀ ਮਦਦ ਨਾਲ ਕਪੂਰ ਨੂੰ ਬਾਲੋ ਅਤੇ ਅੰਬ ਦੀ ਡੰਡੀ ਤੇ ਥੋੜ੍ਹਾ ਜਿਹਾ ਘਿਓ ਪਾ ਦਿਓ।
    \item ਹੌਲੀ ਹੌਲੀ ਸਮੱਗਰੀ ਪਾਓ, ਅੱਗ ਨੂੰ ਬਲਦੀ ਰੱਖਣ ਲਈ ਲੋੜ ਪੈਣ 'ਤੇ ਘਿਓ ਪਾਓ, ਅਤੇ ਅੱਗ ਵਿੱਚ ਪਾਣੀ ਦੀਆਂ ਕੁਝ ਬੂੰਦਾਂ ਛਿੜਕਦੇ ਰਹੋ।
    \item ਵੱਧ ਤੋਂ ਵੱਧ ਲਾਭ ਲੈਣ ਲਈ ਗੁਰਬਾਣੀ ਨੂੰ ਪਿਆਰ ਨਾਲ ਪੜ੍ਹਦੇ ਰਹੋ — ਸ੍ਰੀ ਜਪੁਜੀ ਸਾਹਿਬ, ਸ੍ਰੀ ਜਾਪ ਸਾਹਿਬ, ਸ੍ਰੀ ਅਕਾਲ ਉਸਤਤਿ, ਸ੍ਰੀ ਚੌਪਈ ਸਾਹਿਬ, ਸ੍ਰੀ ਬ੍ਰਹਮ ਕਵਚ, ਜਾਂ ਕੇਵਲ ਮੂਲ ਮੰਤਰ ਅਤੇ ਗੁਰਮੰਤਰ ਦਾ ਜਾਪ।
\end{enumerate}
\textbf{ਨੋਟ:} ਹਵਨ ਦੀਆਂ ਹੋਰ ਵੀ ਕਈ ਕਿਸਮਾਂ ਹਨ, ਜੋ ਪੰਥ ਵਿਚ ਚਲਾਈਆਂ ਜਾਂਦੀਆਂ ਹਨ (ਜੋ ਕਿ ਜ਼ਿਆਦਾਤਰ ਗੁਪਤ ਹਨ)। ਇਹ ਕੇਵਲ ਸਮਾਂ ਬਿਤਾ ਕੇ ਅਤੇ ਕਿਸੇ ਵਿਦਵਾਨ/ਮਹਾਪੁਰਸ਼ ਤੋਂ ਗਿਆਨ ਪ੍ਰਾਪਤ ਕਰਕੇ ਹੀ ਸਿੱਖੇ ਜਾ ਸਕਦੇ ਹਨ। ਇਸ ਲਈ, ਇਹ ਇਸ ਪੁਸਤਕ ਦੇ ਦਾਇਰੇ ਤੋਂ ਬਾਹਰ ਹੈ।

%% ---------------- 9 ----------------
\clearpage
\chapter*{\eng{9.} ਉਦਾਸੀ ਮਹਾਂਪੁਰਖ ਬਾਬਾ ਵਿਰਸਾ ਸਿੰਘ ਜੀ ਮਹਾਰਾਜ ਅਨੁਸਾਰ ਹਵਨ}
\addcontentsline{toc}{chapter}{\eng{9.} ਉਦਾਸੀ ਮਹਾਂਪੁਰਖ ਬਾਬਾ ਵਿਰਸਾ ਸਿੰਘ ਜੀ ਮਹਾਰਾਜ ਅਨੁਸਾਰ ਹਵਨ}
ਭਾਰਤ ਨੂੰ ਹਮੇਸ਼ਾ ਧਰਮ ਦੀ ਧਰਤੀ ਵਜੋਂ ਜਾਣਿਆ ਜਾਂਦਾ ਹੈ। ਰਿਸ਼ੀ ਵੇਦ-ਵਿਆਸ ਜੀ ਮਹਾਰਾਜ ਨੇ ਵੇਦ ਬਾਣੀ (ਰਿਗਵੇਦ, ਯਜੁਰਵੇਦ, ਸਾਮਵੇਦ ਅਤੇ ਅਥਰਵਵੇਦ) ਦੀਆਂ ਚਾਰ ਕਾਪੀਆਂ ਬਣਾਈਆਂ ਸਨ, ਜਿਹਨਾਂ ਵਿਚ ਹਵਨ-ਯੱਗ ਕਰਵਾਉਣ ਦਾ ਜ਼ਿਕਰ ਹੈ। ਗੁਰਬਾਣੀ ਅਨੁਸਾਰ ਵੇਦ ਬਾਣੀ ਗਲਤ ਨਹੀਂ, ਤਾਂ ਹਵਨ ਦਾ ਵਿਰੋਧ ਕਿਉਂ? ਧੰਨ ਸ੍ਰੀ ਗੁਰੂ ਅਰਜਨ ਦੇਵ ਜੀ ਮਹਾਰਾਜ ਦਾ ਹੁਕਮ ਸੀ ਕਿ ਜੇਕਰ ਹਵਨ ਨਾ ਕਰਵਾਇਆ ਜਾਵੇ ਤਾਂ ਅਖੰਡ ਪਾਠ ਅੱਧਾ ਹੀ ਹੋਵੇਗਾ।

ਬਾਬਾ ਵਿਰਸਾ ਸਿੰਘ ਜੀ ਮਹਾਰਾਜ ਫਰਮਾਉਂਦੇ ਹਨ: ਧੰਨ ਸ਼੍ਰੀ ਗੁਰੂ ਗੋਬਿੰਦ ਸਿੰਘ ਜੀ ਮਹਾਰਾਜ ਆਪਣੀ ਬਾਣੀ ਵਿੱਚ ਆਖਦੇ ਹਨ — “ਇਕ ਹੋਮ ਨਾਮੁ ਸੁ ਬੀਰ ॥ ਅਰਿ ਕੀਨ ਜਾਸੁ ਅਧੀਰ ॥” — ਭਾਵ ਹਵਨ ਸੂਰਬੀਰ ਹੈ, ਜੋ ਹਵਨ ਦੀ ਅੱਗ ਵਿਚ ਆਹੂਤੀਆਂ ਚੜ੍ਹਾਉਂਦਾ ਹੈ, ਉਹ ਵੀ ਸੂਰਬੀਰ ਬਣ ਜਾਂਦਾ ਹੈ। ਜੇਕਰ ਕੋਈ ਭਾਵਨਾ ਨਾਲ ਹਵਨ ਕਰਦਾ ਹੈ, ਤਾਂ ਉਹ ਆਪਣੇ ਮਨ ਦੀਆਂ ਇੱਛਾਵਾਂ ਦੀ ਪੂਰਤੀ ਕਰੇਗਾ, ਅਤੇ ਜੇਕਰ ਬ੍ਰਹਮਗਿਆਨੀ ਬਣਨਾ ਚਾਹੁੰਦਾ ਹੈ ਤਾਂ ਉਹ ਯਕੀਨੀ ਤੌਰ ਤੇ ਬਣ ਜਾਵੇਗਾ।

ਪੁਰਾਣੇ ਸਮਿਆਂ ਵਿੱਚ ਰਾਜਿਆਂ ਦੇ ਮਹਿਲਾਂ ਵਿੱਚ ਕਈ ਛੋਟੇ-ਵੱਡੇ ਹਵਨ ਕੁੰਡ ਹੁੰਦੇ ਸਨ, ਜਿਨ੍ਹਾਂ ਵਿੱਚ ਅੱਗ ਬਾਲੀ ਜਾਂਦੀ ਸੀ; ਇਹ ਵਾਯੂਮੰਡਲ ਨੂੰ ਸਾਫ਼ ਰੱਖਣ ਅਤੇ ਉੱਚ-ਵਾਈਬਰੇਸ਼ਨ ਵਾਲੀਆਂ ਰੂਹਾਂ ਨੂੰ ਆਕਰਸ਼ਿਤ ਕਰਨ ਲਈ ਵਰਤਿਆ ਜਾਂਦਾ ਸੀ। ਰਾਜਾ ਜਨਮੇਜਾ ਵੱਲੋਂ ਸ੍ਰੀ ਦਰਬਾਰ ਸਾਹਿਬ, ਅੰਮ੍ਰਿਤਸਰ ਦੀ ਜ਼ਮੀਨ 'ਤੇ ਕਰਵਾਏ ਗਏ ਲੰਬੇ ਹਵਨ ਦਾ ਵੀ ਜ਼ਿਕਰ ਮਿਲਦਾ ਹੈ — ਇਹ ਤਪੱਸਿਆ ਦੀ ਧਰਤੀ ਹੈ ਜੋ ਹੁਣ ਤੀਰਥ ਬਣ ਗਈ ਹੈ।

ਰਾਜਾ ਜਨਕ ਅਤੇ ਰਿਸ਼ੀ ਜਗਵੱਲਕ ਦੀ ਕਥਾ ਅਨੁਸਾਰ — ਹਵਨ ਵਿੱਚ ਸਾਰੇ ਚੜ੍ਹਾਵੇ ਪੂਰੇ ਹੋਣ ਤਾਂ ਚੰਗਾ, ਪਰ ਜੇਕਰ ਪੂਰੀ ਸਮੱਗਰੀ ਨਾ ਮਿਲੇ ਤਾਂ ਕੇਵਲ ਘਿਓ, ਜਾਂ ਫੁੱਲ-ਪੱਤੇ, ਜਾਂ ਕੇਵਲ ਅੱਗ ਨੂੰ ਬਲਦੀ ਰੱਖ ਕੇ ਵਾਹਿਗੁਰੂ ਦੇ ਨਾਮ ਦਾ ਜਾਪ ਕਰਦੇ ਰਹਿਣਾ ਵੀ ਕਾਫ਼ੀ ਹੈ। ਗੁਰਬਾਣੀ ਵਿਚ ਭਗਤ ਰਾਜਾ ਜਨਕ ਦਾ ਜ਼ਿਕਰ ਮਿਲਦਾ ਹੈ: “ਭਗਤੁ ਵਡਾ ਰਾਜਾ ਜਨਕੁ ਹੈ ਗੁਰਮੁਖਿ ਮਾਇਆ ਵਿਚਿ ਉਦਾਸੀ॥”

%% ---------------- 10 ----------------
\clearpage
\chapter*{\eng{10.} ਕੁਝ ਅਕਸਰ ਪੁੱਛੇ ਜਾਣ ਵਾਲੇ ਸਵਾਲ}
\addcontentsline{toc}{chapter}{\eng{10.} ਕੁਝ ਅਕਸਰ ਪੁੱਛੇ ਜਾਣ ਵਾਲੇ ਸਵਾਲ}
\textbf{ਪ੍ਰਸ਼ਨ:} ਕੀ ਹਵਨ, ਹੋਮ ਜਾਂ ਯੱਗ ਇੱਕ ਚੀਜ਼ ਹਨ ਜਾਂ ਵੱਖ-ਵੱਖ?

\textbf{ਉੱਤਰ:} ਹਵਨ ਸੰਸਕ੍ਰਿਤ ਦਾ ਸ਼ਬਦ ਹੈ, ਇਸ ਦਾ ਅਰਥ ਹੋਮ ਜਾਂ ਯੱਗ ਹੈ। ਰਸਮ ਨਿਭਾਉਣ ਵਾਲੇ ਨੂੰ ਅਗਨੀਹੋਤਰੀ ਕਿਹਾ ਜਾਂਦਾ ਹੈ।

\textbf{ਪ੍ਰਸ਼ਨ:} ਇਸ ਰਸਮ ਨਾਲ ਕੀ ਗਲਤ ਹੋ ਸਕਦਾ ਹੈ?

\textbf{ਉੱਤਰ:} ਸ਼ੁਰੂ ਵਿਚ ਪ੍ਰੈਕਟੀਸ਼ਨਰ ਥੋੜ੍ਹਾ ਗ਼ੈਰ-ਹਾਜ਼ਿਰ ਦਿਮਾਗ ਬਣ ਸਕਦਾ ਹੈ ਕਿਉਂਕਿ ਤੁਸੀਂ ਅਜੇ ਵੀ ਸਿੱਖਣ ਦੀ ਪ੍ਰਕਿਰਿਆ ਵਿੱਚ ਹੋ। ਸਮੇਂ ਨਾਲ ਗੱਲ ਥਾਂ 'ਤੇ ਆ ਜਾਵੇਗੀ, ਨਿਰਾਸ਼ ਨਾ ਹੋਵੋ।

\textbf{ਪ੍ਰਸ਼ਨ:} ਤੁਹਾਨੂੰ ਇਹ ਰਸਮ ਕਦੋਂ ਨਹੀਂ ਕਰਨੀ ਚਾਹੀਦੀ?

\textbf{ਉੱਤਰ:} ਜਦੋਂ ਤੁਸੀਂ ਬੁਰੀ ਤਰ੍ਹਾਂ ਬੀਮਾਰ ਹੋ ਅਤੇ ਜਦੋਂ ਤੁਸੀਂ ਇਸ਼ਨਾਨ ਨਹੀਂ ਕੀਤਾ ਹੋਵੇ (ਕੇਸ਼ੀ ਇਸ਼ਨਾਨ ਸਭ ਤੋਂ ਵਧੀਆ ਹੈ)।

\textbf{ਪ੍ਰਸ਼ਨ:} ਕੀ ਮਾਹਵਾਰੀ ਆਉਣ ਵਾਲੀ ਔਰਤ ਹਵਨ ਕਰ ਸਕਦੀ ਹੈ?

\textbf{ਉੱਤਰ:} ਹਵਨ ਊਰਜਾ ਦੀ ਰਸਮ ਹੈ, ਜਿਸ ਵਿੱਚ ਊਰਜਾ ਦਾ ਪ੍ਰਵਾਹ ਹੇਠਲੇ ਚੱਕਰ ਤੋਂ ਉੱਪਰ ਵੱਲ ਹੁੰਦਾ ਹੈ, ਜਦਕਿ ਮਾਹਵਾਰੀ ਦੌਰਾਨ ਇਹ ਪ੍ਰਵਾਹ ਉਲਟਾ ਹੁੰਦਾ ਹੈ। ਇਸ ਲਈ ਇਸ ਸਮੇਂ ਦੌਰਾਨ ਹਵਨ ਦਾ ਅਭਿਆਸ ਕਰਨਾ ਸੁਰੱਖਿਅਤ ਨਹੀਂ ਮੰਨਿਆ ਜਾਂਦਾ।

\textbf{ਪ੍ਰਸ਼ਨ:} ਕੀ ਪਿਛਲੇ ਹਵਨ ਦੀ ਸੁਆਹ ਨੂੰ ਕੁੰਡ ਵਿੱਚੋਂ ਕੱਢਣ ਦੀ ਲੋੜ ਹੈ?

\textbf{ਉੱਤਰ:} ਤਰਜੀਹੀ ਤੌਰ 'ਤੇ ਹਾਂ। ਸੁਆਹ ਨੂੰ ਕੱਚ ਜਾਂ ਤਾਂਬੇ ਦੇ ਭਾਂਡੇ ਵਿੱਚ ਇਕੱਠਾ ਕਰਕੇ ਪੌਦਿਆਂ 'ਤੇ ਛਿੜਕਿਆ ਜਾਂ ਮਿੱਟੀ ਵਿੱਚ ਮਿਲਾਇਆ ਜਾ ਸਕਦਾ ਹੈ।

\textbf{ਪ੍ਰਸ਼ਨ:} ਕੀ ਗਿਆਨ ਪ੍ਰਾਪਤੀ ਲਈ ਹਵਨ ਜ਼ਰੂਰੀ ਹੈ?

\textbf{ਉੱਤਰ:} ਬਿਲਕੁਲ ਜ਼ਰੂਰੀ ਨਹੀਂ। ਹਵਨ ਗੁਰਬਾਣੀ ਅਭਿਆਸ ਲਈ ਮਾਲਾ ਵਰਤਣ ਵਾਂਗ ਹੀ ਇੱਕ ਸਾਧਨ ਹੈ। ਮਨੋਰਥ ਇੱਕ ਹੈ — ਗੁਰਬਾਣੀ ਦਾ ਅਭਿਆਸ। ਜਿਸ ਸਾਧਨ ਨਾਲ ਤੁਸੀਂ ਅਰਾਮਦੇਹ ਮਹਿਸੂਸ ਕਰਦੇ ਹੋ ਉਸ ਦੀ ਵਰਤੋਂ ਕਰੋ।

\textbf{ਪ੍ਰਸ਼ਨ:} ਕੀ ਹਵਨ ਇਸ ਜੀਵਨ ਦੀਆਂ ਮੁਸ਼ਕਲਾਂ ਹੱਲ ਕਰੇਗਾ?

\textbf{ਉੱਤਰ:} ਨਹੀਂ, ਗੁਰਬਾਣੀ ਕਰੇਗੀ। ਹਵਨ ਬਹੁਤ ਸਾਰੇ ਸਾਧਨਾਂ ਵਿੱਚੋਂ ਇੱਕ ਹੈ। ਗੁਰਬਾਣੀ ਦੇ ਅਭਿਆਸ ਨਾਲ ਸੰਚਿਤ ਅਤੇ ਅਗਾਮੀ ਕਰਮ ਬ੍ਰਹਮ ਅਗਨੀ ਨਾਲ ਸੜ ਜਾਂਦੇ ਹਨ, ਜਿਸ ਨਾਲ ਆਤਮਾ ਜਨਮ-ਮਰਨ ਦੇ ਚੱਕਰ ਤੋਂ ਮੁਕਤੀ ਵੱਲ ਵਧਦੀ ਹੈ।

\textbf{ਪ੍ਰਸ਼ਨ:} ਹਵਨ ਦੇ ਸਿਮਰਨ ਨਾਲ ਕੀ ਸੰਬੰਧ ਹੈ?

\textbf{ਉੱਤਰ:} ਹਵਨ ਇੱਕ ਵਾਤਾਵਰਣ ਬਣਾਉਂਦਾ ਹੈ ਜੋ ਮਨ ਨੂੰ ਸ਼ਾਂਤ ਕਰਦਾ ਹੈ, ਇਸ ਤਰ੍ਹਾਂ ਅਭਿਆਸੀ ਨੂੰ ਡੂੰਘੇ ਧਿਆਨ ਲਈ ਤਿਆਰ ਕਰਦਾ ਹੈ। ਗੁਰਮੰਤਰ 'ਤੇ ਧਿਆਨ ਕਰਨ ਨਾਲ ਡੂੰਘੇ ਸਿਮਰਨ ਵਿੱਚ ਮਦਦ ਮਿਲਦੀ ਹੈ।

\textbf{ਪ੍ਰਸ਼ਨ:} ਜੋਤ/ਦੀਵੇ ਦਾ ਕੀ ਮਹੱਤਵ ਹੈ?

\textbf{ਉੱਤਰ:} ਜੋਤ ਨੂੰ ਹਵਨ ਦਾ ਛੋਟਾ ਰੂਪ ਸਮਝਿਆ ਜਾ ਸਕਦਾ ਹੈ। ਘਿਓ ਤੋਂ ਨਿਕਲਣ ਵਾਲੀ ਕੰਪਨ ਊਰਜਾ ਭਲੀ ਹੁੰਦੀ ਹੈ ਅਤੇ ਪ੍ਰਾਰਥਨਾ ਵਾਲੇ ਸਥਾਨ ਨੂੰ ਭਰ ਦਿੰਦੀ ਹੈ, ਗੁਰਬਾਣੀ ਦੇ ਪਾਠ ਦੌਰਾਨ ਇਕਾਗਰਤਾ ਵਿੱਚ ਮਦਦ ਕਰਦੀ ਹੈ।

%% ---------------- 11 ----------------
\clearpage
\chapter*{\eng{11.} ਹਵਾਲੇ}
\addcontentsline{toc}{chapter}{\eng{11.} ਹਵਾਲੇ}
ਇਸ ਪੁਸਤਕ ਨੂੰ ਲਿਖਣ ਵਿੱਚ ਮਦਦ ਕਰਨ ਲਈ ਲੇਖਕ ਸ੍ਰੀ ਗੁਰੂ ਨਾਨਕ ਦੇਵ ਜੀ ਮਹਾਰਾਜ ਅਤੇ ਬਾਬਾ ਸ੍ਰੀ ਚੰਦ ਜੀ ਮਹਾਰਾਜ ਦਾ ਤਹਿ ਦਿਲੋਂ ਧੰਨਵਾਦ ਕਰਦਾ ਹੈ। ਜੇ ਪਾਠਕ ਨੂੰ ਕੋਈ ਗਲਤੀ ਮਿਲਦੀ ਹੈ, ਤਾਂ ਉਹ ਇਕੱਲੀ ਲੇਖਕ ਦੀ ਹੈ।
\begin{enumerate}[label=\arabic*.]
    \item ਸ੍ਰੀ ਗੁਰੂ ਗ੍ਰੰਥ ਸਾਹਿਬ ਜੀ
    \item ਸ੍ਰੀ ਦਸਮ ਗ੍ਰੰਥ ਸਾਹਿਬ ਜੀ
    \item ਸ੍ਰੀ ਸਰਬਲੋਹ ਗ੍ਰੰਥ ਸਾਹਿਬ ਜੀ
    \item ਭਾਈ ਗੁਰਦਾਸ ਜੀ ਦੀਆਂ ਵਾਰਾਂ
    \item ਸੂਰਜ ਪ੍ਰਕਾਸ਼ ਗ੍ਰੰਥ
    \item ਪੰਥ ਪ੍ਰਕਾਸ਼ ਗ੍ਰੰਥ
    \item ਮੁਕਤੀ ਮਾਰਗ ਗ੍ਰੰਥ
    \item ਸੰਨਿਆ ਗਿਆਨਸ਼ਿਰਤੀ ਅਤੇ ਸ਼੍ਰੀ ਵਿਧਿਆਨੰਦ ਦੁਆਰਾ ਯਗਿਆ — ਇੱਕ ਵਿਆਪਕ ਸਮੀਖਿਆ
    \item ਏਲੀ ਹੋਡਲ ਦੁਆਰਾ ਧਰਤੀ ਉੱਤੇ ਅਗਨੀਹੋਤਰ ਹਵਨ
    \item ਗੁਰੂ ਖਾਲਸਾ ਪੰਥ ਦੇ ਬਜ਼ੁਰਗ ਨਿਹੰਗ ਸਿੰਘਾਂ ਦੇ ਮੌਖਿਕ ਸਰੋਤ
    \item ਬਜ਼ੁਰਗ ਨਾਮਧਾਰੀ ਸਿੱਖ
    \item ਵੇਦ
    \item ਪੁਰਾਣ
\end{enumerate}

%% ---------------- ABOUT THE AUTHOR ----------------
\clearpage
\chapter*{ਲੇਖਕ ਬਾਰੇ}
\addcontentsline{toc}{chapter}{ਲੇਖਕ ਬਾਰੇ}
ਭਾਈ ਮਨਪ੍ਰੀਤ ਸਿੰਘ ਜੀ, ਨਾਨਕ ਨਿਵਾਸ — ਗੁਰਮਤਿ ਵਿਦਿਆਲਾ ਦੇ ਸੰਸਥਾਪਕ, ਲਗਭਗ ਅੱਠ ਸਾਲਾਂ ਤੋਂ ਬ੍ਰਹਮ ਵਿਦਿਆ ਦੀ ਸਿੱਖਿਆ ਦੇ ਰਹੇ ਹਨ, ਅਤੇ ਉਦਾਸੀ, ਨਿਰਮਲਾ ਅਤੇ ਨਿਹੰਗ ਪਰੰਪਰਾਵਾਂ ਵਿੱਚ ਵਿਰਾਸਤੀ ਜੜ੍ਹਾਂ ਰੱਖਦੇ ਹਨ।

%% ---------------- PHOTO GALLERY ----------------
\newpage
\thispagestyle{empty}
\vspace*{\fill}

\begin{center}
    {\Huge \textbf{ਚਿੱਤਰ ਸੰਗ੍ਰਹਿ}}
    \addcontentsline{toc}{chapter}{ਚਿੱਤਰ ਸੰਗ੍ਰਹਿ}
\end{center}

\vspace*{\fill}
\galleryphoto{katha_prakash.jpeg}{ਸ੍ਰੀ ਗੁਰੂ ਗ੍ਰੰਥ ਸਾਹਿਬ, ਸ੍ਰੀ ਦਸਮ ਗ੍ਰੰਥ ਸਾਹਿਬ ਅਤੇ ਸ੍ਰੀ ਸਰਬਲੋਹ ਗ੍ਰੰਥ ਸਾਹਿਬ ਜੀ ਦਾ ਕਥਾ ਪ੍ਰਕਾਸ਼ (ਸ਼ਹੀਦੀ ਬਾਗ਼, ਸ੍ਰੀ ਅਨੰਦਪੁਰ ਸਾਹਿਬ ਵਿਖੇ)}
\galleryphoto{giani_sher_singh.jpg}{ਮਹਾਕਾਲ ਬਾਬਾ ਗਿਆਨੀ ਸ਼ੇਰ ਸਿੰਘ ਜੀ}
\galleryphoto{giani_sukha_singh_ji.jpeg}{ਮਹਾਕਾਲ ਬਾਬਾ ਸੁੱਖਾ ਸਿੰਘ ਜੀ}
\galleryphoto{satguru_ram_singh_ji.jpg}{ਸ਼੍ਰੀ ਸਤਿਗੁਰੂ ਰਾਮ ਸਿੰਘ ਜੀ}
\galleryphoto{baba_shri_chand_ji.jpg}{ਬਾਬਾ ਸ਼੍ਰੀ ਚੰਦ ਜੀ}
\galleryphoto{sant_darshan_singh_ji_sashtri.png}{ਸੰਤ ਦਰਸ਼ਨ ਸਿੰਘ ਜੀ ਸ਼ਾਸਤਰੀ}
\galleryphoto{giani_gurvinder_singh_nangli.jpeg}{ਗਿਆਨੀ ਗੁਰਵਿੰਦਰ ਸਿੰਘ ਜੀ ਮੁਖ ਸੇਵਾਦਾਰ ਸ੍ਰੀ ਸਰਬਲੋਹ ਬੁੰਗਾ ਨੰਗਲੀ ਸਾਹਿਬ}
\galleryphoto{baba_nand_singh_ji.jpeg}{ਬਾਬਾ ਨੰਦ ਸਿੰਘ ਜੀ}
\galleryphoto{bhai_manpreet_singh_ji.jpeg}{ਭਾਈ ਮਨਪ੍ਰੀਤ ਸਿੰਘ}
}% end \punjabifont
\end{document}
